\documentclass[11pt,a4paper]{article}
\usepackage{fontspec,amsmath,amssymb,amsthm,microtype}
\setmainfont{lmroman10-regular.otf}[BoldFont=lmroman10-bold.otf,ItalicFont=lmroman10-italic.otf,BoldItalicFont=lmroman10-bolditalic.otf]
\usepackage{xcolor}
\usepackage[most]{tcolorbox}
\usepackage{mathrsfs}
\providecommand{\lt}{<}
\providecommand{\gt}{>}
\setmonofont{DejaVuSansMono.ttf}[BoldFont=DejaVuSansMono-Bold.ttf,ItalicFont=DejaVuSansMono-Oblique.ttf,BoldItalicFont=DejaVuSansMono-BoldOblique.ttf,Scale=0.86]
\newfontfamily\leanfallback{STIXTwoMath-Regular.otf}[Scale=0.86]
\definecolor{leanrule}{gray}{0.62}
\definecolor{leanback}{gray}{0.962}
\definecolor{leantitle}{gray}{0.22}
\definecolor{leanmuted}{gray}{0.38}
\newtcolorbox{leanbox}[1]{enhanced,breakable,colback=leanback,colframe=leanrule,
  boxrule=0.4pt,arc=0.6mm,left=2.4mm,right=2.4mm,top=1.8mm,bottom=1.8mm,
  before skip=0.8em,after skip=1em,fonttitle=\small\bfseries,coltitle=leantitle,
  colbacktitle=leanback,titlerule=0.3pt,title={#1}}
\newlength{\leancharwidth}
\newenvironment{leancode}{\par\smallskip\begingroup\ttfamily\footnotesize\raggedright
  \settowidth{\leancharwidth}{M}\setlength{\parskip}{0pt}\setlength{\parindent}{0pt}
  \frenchspacing}{\par\endgroup\smallskip}
\newcommand{\leanline}[2]{\hangindent=\dimexpr#1\leancharwidth+4\leancharwidth\relax
  \hangafter=1\noindent\hspace*{#1\leancharwidth}#2\par}
\newcommand{\leansource}[1]{\par{\footnotesize\color{leanmuted}#1}\par}
\newcommand{\leanlabel}[1]{\par\smallskip\noindent{\small\bfseries #1}\par}
\usepackage[margin=25mm]{geometry}
\usepackage[unicode,hidelinks]{hyperref}
\hypersetup{pdfauthor={LangevinLean Project}}
\setlength{\emergencystretch}{3em}
\setlength{\parskip}{0.4em}
\setlength{\parindent}{0pt}
\allowdisplaybreaks[2]
\raggedbottom
\title{\texorpdfstring{Kinetic Interacting Particle Langevin Monte Carlo}{Kinetic Interacting Particle Langevin Monte Carlo}}
\author{LangevinLean Project}
\date{}
\begin{document}\maketitle
Draft mathematical exposition; these notes may contain mistakes.
The Brownian KIPLD convergence proposition and both Gaussian algorithm parameter-error theorems are proved. KIPLMC1 reuses the shared exponential integrator; KIPLMC2 has an internal OBABO contraction and Gibbs sampling proof. The source audit records the exact correspondence and source corrections.
\section{\texorpdfstring{Theorem}{Theorem}}


Let \(N\ge1\), and suppose the joint potential \(U(\theta,x)\) is differentiable, \(m\)-strongly convex, and has \(L\)-Lipschitz gradient, with \(0<m\le L\). Initialize the scaled position \(Z_0\sim\nu\) with finite second moment. Independently initialize scaled velocity as \(\mathcal N(0,N^{-1}I)\), and use the independent Brownian source. If \(\gamma^2\ge m+L\), the constructed physical KIPLD parameter law satisfies


\[\boxed{W_2(\mathcal L(\Theta_t),\pi_\Theta)\le\sqrt2e^{-mt/\gamma}\bigl(\mathbb E\|Z_0-\bar Z_*\|^2\bigr)^{1/2}.}\]


Here \(\bar Z_*\sim\pi_Z\) is independent of \(Z_0\), and \(\pi_Z(dz)\propto e^{-V_N(z)}dz\). The target parameter marginal is exactly \(\pi_\Theta(d\theta)\propto\ell(\theta)^N d\theta\). This is the supplied manuscript’s continuous convergence proposition.




The proof applies the earlier kinetic Gibbs-mixing theorem (see \cite[Section~3.58.4]{LangevinLean2026Kinetic}) with \(c=\sqrt N\), unit-temperature friction \(c\gamma\), and time \(t/c\). The sum potential has curvature \(Nm\), so its mixing exponent becomes \((Nm)/(c\gamma)\,(t/c)=mt/\gamma\). The flow is rescaled in velocity only; its position endpoint is unchanged. A separate theorem verifies its full physical integral equation and its independent thermal initial law. The earlier Brownian clock-scaling theorem proves that the transformed forcing has amplitude \(\sqrt{2\gamma/N}\) (see \cite[Section~3.105.1]{LangevinLean2026Kinetic}, \cite[Section~3.105.2]{LangevinLean2026Kinetic}). Projection to the parameter cannot increase a coupling’s cost, and the independent initial coupling gives the displayed constant.



This statement bounds transport to the powered parameter marginal, rather than distance to its maximizer. To obtain parameter RMS error around the maximizer, one subsequently adds the sharp equilibrium concentration radius from the earlier proximal report (see \cite[Section~3.6.1]{LangevinLean2026Proximal}). Time reduces the displayed transport term, while particle count controls that additional statistical radius.

\begin{leanbox}{Lean: \texttt{kipld\_\allowbreak{}parameter\_\allowbreak{}sampling\_\allowbreak{}bound}}
\leansource{\texttt{kipld\_\allowbreak{}parameter\_\allowbreak{}sampling\_\allowbreak{}bound} \quad KIPLMC/\allowbreak{}Continuous.lean, lines 222–245}
\begin{leancode}
\leanline{0}{theorem\ kipld\_parameter\_sampling\_bound\ \{N\ dθ\ dx\ :\ ℕ\}\ (hN\ :\ 0<N)}
\leanline{4}{(U\ :\ IPLA.Joint\ dθ\ dx\ →\ ℝ)\ (m\ :\ ℝ)\ (L\ γ\ :\ ℝ≥0)}
\leanline{4}{(hdiff\ :\ Differentiable\ ℝ\ U)\ (hsc\ :\ StrongConvexOn\ univ\ m\ U)}
\leanline{4}{(hLip\ :\ LipschitzWith\ L\ (gradient\ U))\ (hmL\ :\ m≤(L:ℝ))\ (hm\ :\ 0<m)}
\leanline{4}{(hγ\ :\ 0<(γ:ℝ))\ (hfriction\ :\ m+(L:ℝ)≤(γ:ℝ)\textasciicircum{}2)}
\leanline{4}{(ν\ :\ Measure\ (State\ N\ dθ\ dx))\ [IsProbabilityMeasure\ ν]}
\leanline{4}{(hν\ :\ MemLp\ (id\ :\ State\ N\ dθ\ dx\ →\ State\ N\ dθ\ dx)\ 2\ ν)\ (T\ :\ ℝ)\ (hT\ :\ 0≤T)\ :}
\leanline{4}{ULA.wasserstein2Real}
\leanline{6}{((kipldBrownianPositionLaw\ hN\ U\ L\ γ\ hdiff\ hLip\ ν\ hν\ T\ hT).map\ IPLA.parameterCLM)}
\leanline{6}{(IPLA.marginalTarget\ N\ U)\ ≤\ Real.sqrt\ 2*Real.exp\ ((-m/(γ:ℝ))*T)*}
\leanline{8}{Real.sqrt\ (∫\ p\ :\ State\ N\ dθ\ dx×State\ N\ dθ\ dx,\ ‖p.1-p.2‖\textasciicircum{}2\ ∂ν.prod\ (IPLA.gibbsTarget\ U))}
\end{leancode}
\end{leanbox}

\section{\texorpdfstring{The likelihood and interacting latent particles}{The likelihood and interacting latent particles}}


The observation \(y\) is fixed. A model gives a joint density
\(p_\theta(x,y)\), but the quantity to optimize is its integral over an
unobserved variable:


\[p_\theta(y)=\int p_\theta(x,y)dx,
\qquad\theta_*=\arg\max_\theta\log p_\theta(y).\]


Computing this integral and its gradient exactly can be difficult. The
particle method retains \(N\) latent states \(X^1,\ldots,X^N\), coupled
through the same parameter \(\theta\). The force on the parameter is the
average of the parameter gradients, while the latent drift uses the corresponding
latent gradient:


\[F^\theta=\frac1N\sum_{i=1}^N\nabla_\theta U(\theta,X^i),
\qquad F^{x_i}=\nabla_xU(\theta,X^i).\]


Adding a parameter noise of size \(N^{-1/2}\) makes the whole system a
Langevin system with an explicit Gibbs equilibrium. The particle number then
has two roles: it provides interacting latent samples and determines how
strongly the stationary parameter distribution concentrates near the optimizer.



The common parameter is represented once while each latent variable has its own coordinate. Conditional on that parameter, the equilibrium latent integrations factor into identical marginal likelihood integrals. Increasing particle count changes the parameter density by increasing the power of the likelihood, and increases the state dimension linearly.

\begin{leanbox}{Lean: \texttt{existsUnique\_\allowbreak{}mmle}, \texttt{gibbsPhaseTarget\_\allowbreak{}parameter}}
\leansource{\texttt{existsUnique\_\allowbreak{}mmle} \quad KIPLMC/\allowbreak{}Target.lean, lines 65–69}
\begin{leancode}
\leanline{0}{theorem\ existsUnique\_mmle\ \{dθ\ dx\ :\ ℕ\}\ (U\ :\ IPLA.Joint\ dθ\ dx\ →\ ℝ)}
\leanline{4}{(m\ :\ ℝ)\ (hm\ :\ 0\ <\ m)\ (hdiff\ :\ Differentiable\ ℝ\ U)}
\leanline{4}{(hsc\ :\ StrongConvexOn\ univ\ m\ U)\ :}
\leanline{4}{∃!\ θ,\ IPLA.IsMarginalMaximizer\ U\ θ}
\end{leancode}
\leansource{\texttt{gibbsPhaseTarget\_\allowbreak{}parameter} \quad KIPLMC/\allowbreak{}Target.lean, lines 54–62}
\begin{leancode}
\leanline{0}{theorem\ gibbsPhaseTarget\_parameter\ \{N\ dθ\ dx\ :\ ℕ\}\ (hN\ :\ 0\ <\ N)}
\leanline{4}{(U\ :\ IPLA.Joint\ dθ\ dx\ →\ ℝ)\ (m\ :\ ℝ)\ (hm\ :\ 0\ <\ m)}
\leanline{4}{(hdiff\ :\ Differentiable\ ℝ\ U)\ (hsc\ :\ StrongConvexOn\ univ\ m\ U)\ :}
\leanline{4}{(@gibbsPhaseTarget\ N\ dθ\ dx\ U).map\ (fun\ z\ ↦\ IPLA.parameterCLM\ z.2)\ =}
\leanline{6}{IPLA.marginalTarget\ N\ U}
\end{leancode}
\leanlabel{How the Lean proof follows the argument}
\texttt{exists\allowbreak{}Unique\_\allowbreak{}mmle} supplies the existence and uniqueness of the marginal maximizer under the stated joint assumptions. \texttt{gibbs\allowbreak{}Phase\allowbreak{}Target\_\allowbreak{}parameter} identifies the parameter projection of the normalized phase target with the powered marginal. The proof uses the normalized measure, so the probabilistic parameter error is measured against an actual probability distribution.
\end{leanbox}

\section{\texorpdfstring{Joint strong convexity and smoothness of the statistical potential}{Joint strong convexity and smoothness of the statistical potential}}


The source paper assumes \(U\in C^2\), \(\mu>0\), \(L>0\), and


\[\langle z-z',\nabla U(z)-\nabla U(z')\rangle
\ge\mu\|z-z'\|^2,\qquad
\|\nabla U(z)-\nabla U(z')\|\le L\|z-z'\|,\]


Here \(z=(\theta,x)\). The proved convergence bounds require differentiability and the displayed gradient inequalities. Under the source’s additional twice-continuous-differentiability assumption, these inequalities are equivalent to
\(\mu I\preceq\nabla^2U(z)\preceq LI\). Thus \(\mu\le L\) in positive
dimension. Strong convexity in \(x\) with \(\theta\) fixed does not imply joint strong convexity: the lower bound must also hold in directions that change both variables.




The paper's unregularized logistic example illustrates the distinction.
Its Gaussian prior contributes a quadratic in \(x-\theta\), which leaves
common shifts partly uncontrolled; logistic curvature can vanish far from the
origin. Spanning covariates give strict convexity, while a uniform positive
strong-convexity constant requires additional regularization. The source itself
separates these experiments from its strong-convexity theorem.



The required lower curvature concerns arbitrary joint directions \((u,w)\), including directions that move parameter and latent variables together. Positive curvature in either coordinate with the other fixed is insufficient. A direction in which both move can remove a prior penalty and expose weak or vanishing curvature in the likelihood.

\begin{leanbox}{Lean: \texttt{averagePotential\_\allowbreak{}strongConvex}, \texttt{force\_\allowbreak{}lipschitz}, \texttt{jointPotential\_\allowbreak{}gradient\_\allowbreak{}lipschitz} and 3 more}
\leansource{\texttt{averagePotential\_\allowbreak{}strongConvex} \quad KIPLMC/\allowbreak{}Geometry.lean, lines 48–58}
\begin{leancode}
\leanline{0}{theorem\ averagePotential\_strongConvex\ \{N\ dθ\ dx\ :\ ℕ\}\ (hN\ :\ 0\ <\ N)}
\leanline{4}{\{U\ :\ IPLA.Joint\ dθ\ dx\ →\ ℝ\}\ \{m\ :\ ℝ\}}
\leanline{4}{(hU\ :\ StrongConvexOn\ Set.univ\ m\ U)\ :}
\leanline{4}{StrongConvexOn\ Set.univ\ m\ (@averagePotential\ N\ dθ\ dx\ U)}
\end{leancode}
\leansource{\texttt{force\_\allowbreak{}lipschitz} \quad KIPLMC/\allowbreak{}Geometry.lean, lines 82–98}
\begin{leancode}
\leanline{0}{theorem\ force\_lipschitz\ \{N\ dθ\ dx\ :\ ℕ\}\ (hN\ :\ 0\ <\ N)}
\leanline{4}{\{U\ :\ IPLA.Joint\ dθ\ dx\ →\ ℝ\}\ \{L\ :\ ℝ≥0\}}
\leanline{4}{(hdiff\ :\ Differentiable\ ℝ\ U)\ (hLip\ :\ LipschitzWith\ L\ (gradient\ U))\ :}
\leanline{4}{LipschitzWith\ L\ (@force\ N\ dθ\ dx\ U)}
\end{leancode}
\leansource{\texttt{jointPotential\_\allowbreak{}gradient\_\allowbreak{}lipschitz} \quad KIPLMC/\allowbreak{}LogisticExample.lean, lines 322–330}
\begin{leancode}
\leanline{0}{theorem\ jointPotential\_gradient\_lipschitz\ [CompleteSpace\ E]}
\leanline{4}{(tau\ :\ ℝ≥0)\ (v\ :\ J→E)\ (y\ :\ J→ℝ)\ :}
\leanline{4}{LipschitzWith\ (2*tau+∑\ j,‖v\ j‖₊\textasciicircum{}2/4)}
\leanline{6}{(gradient\ (jointPotential\ (tau:ℝ)\ v\ y))}
\end{leancode}
\leansource{\texttt{jointHessian\_\allowbreak{}quadratic} \quad KIPLMC/\allowbreak{}LogisticExample.lean, lines 433–448}
\begin{leancode}
\leanline{0}{theorem\ jointHessian\_quadratic\ (tau\ :\ ℝ)\ (v\ :\ J→E)\ (z\ d\ :\ Joint\ E)\ :}
\leanline{4}{{\leanfallback ⟪}jointHessian\ tau\ v\ z\ d,d{\leanfallback ⟫}=}
\leanline{6}{curvature\ tau\ v\ (WithLp.ofLp\ z).2\ (WithLp.ofLp\ d).1\ (WithLp.ofLp\ d).2}
\end{leancode}
\leansource{\texttt{potential\_\allowbreak{}strictConvex\_\allowbreak{}of\_\allowbreak{}span} \quad KIPLMC/\allowbreak{}LogisticExample.lean, lines 172–195}
\begin{leancode}
\leanline{0}{theorem\ potential\_strictConvex\_of\_span\ (tau\ :\ ℝ)\ (htau\ :\ 0<tau)\ (v\ :\ J→E)\ (y\ :\ J→ℝ)}
\leanline{4}{(hspan\ :\ Submodule.span\ ℝ\ (Set.range\ v)=⊤)\ :}
\leanline{4}{StrictConvexOn\ ℝ\ Set.univ\ (potential\ tau\ v\ y)}
\end{leancode}
\leansource{\texttt{jointPotential\_\allowbreak{}not\_\allowbreak{}strongConvex} \quad KIPLMC/\allowbreak{}LogisticExample.lean, lines 478–512}
\begin{leancode}
\leanline{0}{theorem\ jointPotential\_not\_strongConvex\ (tau\ :\ ℝ)\ (v\ :\ J→E)\ (y\ :\ J→ℝ)}
\leanline{4}{(w\ :\ E)\ (hw\ :\ w≠0)\ (m\ :\ ℝ)\ (hm\ :\ 0<m)\ :}
\leanline{4}{¬StrongConvexOn\ Set.univ\ m\ (jointPotential\ tau\ v\ y)}
\end{leancode}
\leanlabel{How the Lean proof follows the argument}
\texttt{average\allowbreak{}Potential\_\allowbreak{}strong\allowbreak{}Convex} divides the lifted sum-potential curvature by the positive particle count. \texttt{force\_\allowbreak{}lipschitz} performs the same normalization for the gradient Lipschitz constant. Both use the shared particle geometry, whose squared joint-coordinate norms sum to the particle count times the scaled state norm. Positivity of particle count is needed for cancelling this factor.
\end{leanbox}

\section{\texorpdfstring{The kinetic particle equations in physical coordinates}{The kinetic particle equations in physical coordinates}}


KIPLD adds velocity blocks \(V^\theta,V^{x_i}\) and uses the same
friction in each:


\[d\theta=V^\theta dt,\qquad dX^i=V^{x_i}dt,\]
\[dV^\theta=-(\gamma V^\theta+F^\theta)dt
+\sqrt{2\gamma/N}\,dB^0,
\quad dV^{x_i}=-(\gamma V^{x_i}+F^{x_i})dt+\sqrt{2\gamma}\,dB^i.\]


The Brownian blocks are independent. The parameter velocity has equilibrium
covariance \(N^{-1}I\); every latent velocity has covariance \(I\). These
scales are chosen together with the average parameter force. Changing one
without changing the others generally changes the equilibrium target.




The noise acts on velocity, while positions are integrals of velocities.
This creates the smoother position paths responsible for the exponential
integrator's improved discretization dependence. The target remains a
statistical estimation problem: the final parameter distribution should
concentrate around \(\theta_*\) as the particle number grows.



The parameter noise is reduced by the square root of particle count, matching its averaged force and concentrated equilibrium marginal. Each physical latent equation retains unit temperature. Scaling all latent positions and velocities by \(N^{-1/2}\) turns these different physical amplitudes into one uniform phase-space temperature.

\begin{leanbox}{Lean: \texttt{kipldBrownianFlow\_\allowbreak{}integral}, \texttt{kipldBrownianFlow\_\allowbreak{}initial\_\allowbreak{}hasLaw}}
\leansource{\texttt{kipldBrownianFlow\_\allowbreak{}integral} \quad KIPLMC/\allowbreak{}Continuous.lean, lines 188–219}
\begin{leancode}
\leanline{0}{theorem\ kipldBrownianFlow\_integral\ \{N\ dθ\ dx\ :\ ℕ\}\ (hN\ :\ 0<N)}
\leanline{4}{(U\ :\ IPLA.Joint\ dθ\ dx\ →\ ℝ)\ (L\ γ\ :\ ℝ≥0)}
\leanline{4}{(hdiff\ :\ Differentiable\ ℝ\ U)\ (hLip\ :\ LipschitzWith\ L\ (gradient\ U))}
\leanline{4}{(ν\ :\ Measure\ (State\ N\ dθ\ dx))\ [IsProbabilityMeasure\ ν]}
\leanline{4}{(hν\ :\ MemLp\ (id\ :\ State\ N\ dθ\ dx\ →\ State\ N\ dθ\ dx)\ 2\ ν)\ (T\ :\ ℝ)\ (hT\ :\ 0≤T)\ :}
\leanline{4}{let\ P\ :=\ ν.prod\ (kipldBrownianSource\ N\ dθ\ dx)}
\leanline{4}{let\ S\ :=\ kipldBrownianFlow\ hN\ U\ L\ γ\ hdiff\ hLip\ ν\ hν\ T\ hT}
\leanline{4}{let\ B\ :=\ ProductContinuous.input\ (kipldBrownianSource\ N\ dθ\ dx)\ ν}
\leanline{6}{(CanonicalContinuous.input\ (Coordinate\ N\ dθ\ dx)\ (stdGaussian\ (State\ N\ dθ\ dx))\ (kipldClockFriction\ N\ γ))}
\leanline{4}{Continuous\ S\ ∧\ ∀\ t\ ∈\ Icc\ 0\ T,\ S\ t-S\ 0\ =}
\leanline{6}{-(∫\ s\ in\ 0..t,kineticNegativeDrift\ (ULA.lpDrift\ P\ (force\ U)\ (force\_lipschitz\ hN\ hdiff\ hLip))\ γ\ 1\ (S\ s))+}
\leanline{6}{((Real.sqrt\ N)⁻¹\ •\ B\ (t/Real.sqrt\ N),0)}
\end{leancode}
\leansource{\texttt{kipldBrownianFlow\_\allowbreak{}initial\_\allowbreak{}hasLaw} \quad KIPLMC/\allowbreak{}Continuous.lean, lines 58–95}
\begin{leancode}
\leanline{0}{theorem\ kipldBrownianFlow\_initial\_hasLaw\ \{N\ dθ\ dx\ :\ ℕ\}\ (hN\ :\ 0<N)}
\leanline{4}{(U\ :\ IPLA.Joint\ dθ\ dx\ →\ ℝ)\ (L\ γ\ :\ ℝ≥0)}
\leanline{4}{(hdiff\ :\ Differentiable\ ℝ\ U)\ (hLip\ :\ LipschitzWith\ L\ (gradient\ U))}
\leanline{4}{(ν\ :\ Measure\ (State\ N\ dθ\ dx))\ [IsProbabilityMeasure\ ν]}
\leanline{4}{(hν\ :\ MemLp\ (id\ :\ State\ N\ dθ\ dx\ →\ State\ N\ dθ\ dx)\ 2\ ν)\ (T\ :\ ℝ)\ (hT\ :\ 0≤T)\ :}
\leanline{4}{HasLaw\ (lpPhaseCoe\ (ν.prod\ (kipldBrownianSource\ N\ dθ\ dx))}
\leanline{6}{(kipldBrownianFlow\ hN\ U\ L\ γ\ hdiff\ hLip\ ν\ hν\ T\ hT\ 0))}
\leanline{6}{((velocityTarget\ N\ dθ\ dx).prod\ ν)\ (ν.prod\ (kipldBrownianSource\ N\ dθ\ dx))}
\end{leancode}
\leanlabel{How the Lean proof follows the argument}
\texttt{kipld\allowbreak{}Brownian\allowbreak{}Flow\_\allowbreak{}integral} is the physical-time integral equation of the constructed particle flow. \texttt{kipld\allowbreak{}Brownian\allowbreak{}Flow\_\allowbreak{}initial\_\allowbreak{}has\allowbreak{}Law} verifies the initial thermal phase law on its product source. The position and velocity coordinates in the integral equation are therefore those of the same random process used later in the transport bound.
\end{leanbox}

\section{\texorpdfstring{Transformation to unit-temperature kinetic coordinates}{Transformation to unit-temperature kinetic coordinates}}


The particle system looks anisotropic, but it becomes a standard kinetic
system in \(d_z=d_\theta+Nd_x\) dimensions. At rescaled time \(t\), define


\[\widetilde Z_t=(\theta_{\sqrt Nt},N^{-1/2}X^1_{\sqrt Nt},\ldots),
\qquad\widetilde V_t=(\sqrt N V^\theta_{\sqrt Nt},V^{x_1}_{\sqrt Nt},\ldots).\]


Set


\[\bar U_N(\theta,z_1,\ldots,z_N)=\frac1N\sum_{i=1}^NU(\theta,\sqrt N z_i),
\qquad f_N=N\bar U_N,\qquad\widetilde\gamma=\sqrt N\gamma.\]


Brownian scaling produces independent standard blocks
\(\widetilde B_t^i=N^{-1/4}B^i_{\sqrt Nt}\). Direct substitution yields


\[d\widetilde Z_t=\widetilde V_tdt,\qquad
d\widetilde V_t=-(\widetilde\gamma\widetilde V_t+\nabla f_N(\widetilde Z_t))dt
+\sqrt{2\widetilde\gamma}\,d\widetilde B_t.\]


This is Proposition A.1. The parameter position stays unchanged under the
spatial scaling, which will make projection of the final transport estimate
particularly simple. An original step \(\eta\) corresponds to a rescaled
step \(\widetilde\eta=\eta/\sqrt N\).



The latent spatial scaling is applied first. The subsequent velocity and clock transformation fixes these scaled position coordinates. The parameter position is unchanged throughout. Friction and step are changed inversely, preserving their dimensionless product. The normalized Gaussian velocity and transformed Brownian amplitude must be transformed together; a deterministic clock identity alone would not identify the sampler’s probability law.

\begin{leanbox}{Lean: \texttt{kinetic\_\allowbreak{}integral\_\allowbreak{}clock}, \texttt{clockBrownian\_\allowbreak{}isPreBrownian}, \texttt{kipldBrownianPositionLaw\_\allowbreak{}eq\_\allowbreak{}canonical}}
\leansource{\texttt{kinetic\_\allowbreak{}integral\_\allowbreak{}clock} \quad KIPLMC/\allowbreak{}KineticClock.lean, lines 29–57}
\begin{leancode}
\leanline{0}{theorem\ kinetic\_integral\_clock\ (g\ :\ E\ →\ E)\ (hg\ :\ Continuous\ g)\ (c\ γ\ T\ :\ ℝ)\ (hc\ :\ 0<c)}
\leanline{4}{(initial\ :\ E×E)\ (B\ :\ ℝ\ →\ E)\ (Y\ :\ ℝ\ →\ E×E)\ (hY\ :\ Continuous\ Y)}
\leanline{4}{(heY\ :\ ∀\ s\ ∈\ Icc\ 0\ (T/c),Y\ s-initial\ =}
\leanline{6}{-(∫\ r\ in\ 0..s,KineticLangevin.kineticNegativeDrift\ g\ (c*γ)\ 1\ (Y\ r))+(B\ s,0))\ :}
\leanline{4}{∀\ t\ ∈\ Icc\ 0\ T,}
\leanline{6}{rescaleVelocity\ c⁻¹\ (Y\ (t/c))-rescaleVelocity\ c⁻¹\ initial\ =}
\leanline{8}{-(∫\ s\ in\ 0..t,KineticLangevin.kineticNegativeDrift\ (fun\ x\ ↦\ (c\textasciicircum{}2)⁻¹\ •\ g\ x)\ γ\ 1}
\leanline{10}{(rescaleVelocity\ c⁻¹\ (Y\ (s/c))))+(c⁻¹\ •\ B\ (t/c),0)}
\end{leancode}
\leansource{\texttt{clockBrownian\_\allowbreak{}isPreBrownian} \quad KineticLangevin/\allowbreak{}BrownianScaling.lean, lines 15–19}
\begin{leancode}
\leanline{0}{theorem\ clockBrownian\_isPreBrownian\ \{ι\ :\ Type*\}\ (c\ :\ ℝ≥0)\ (hc\ :\ c≠0)\ (i\ :\ ι)\ :}
\leanline{4}{IsPreBrownianReal\ (clockBrownian\ c\ i)\ (ULA.CanonicalBrownian.vectorProbability\ ι)}
\end{leancode}
\leansource{\texttt{kipldBrownianPositionLaw\_\allowbreak{}eq\_\allowbreak{}canonical} \quad KIPLMC/\allowbreak{}Continuous.lean, lines 45–55}
\begin{leancode}
\leanline{0}{theorem\ kipldBrownianPositionLaw\_eq\_canonical\ \{N\ dθ\ dx\ :\ ℕ\}\ (hN\ :\ 0<N)}
\leanline{4}{(U\ :\ IPLA.Joint\ dθ\ dx\ →\ ℝ)\ (L\ γ\ :\ ℝ≥0)}
\leanline{4}{(hdiff\ :\ Differentiable\ ℝ\ U)\ (hLip\ :\ LipschitzWith\ L\ (gradient\ U))}
\leanline{4}{(ν\ :\ Measure\ (State\ N\ dθ\ dx))\ [IsProbabilityMeasure\ ν]}
\leanline{4}{(hν\ :\ MemLp\ (id\ :\ State\ N\ dθ\ dx\ →\ State\ N\ dθ\ dx)\ 2\ ν)}
\leanline{4}{(T\ :\ ℝ)\ (hT\ :\ 0≤T)\ :}
\leanline{4}{kipldBrownianPositionLaw\ hN\ U\ L\ γ\ hdiff\ hLip\ ν\ hν\ T\ hT\ =}
\leanline{6}{CanonicalContinuous.positionLaw\ (Coordinate\ N\ dθ\ dx)\ (stdGaussian\ (State\ N\ dθ\ dx))}
\leanline{8}{IsGaussian.memLp\_two\_id\ ν\ hν\ (@IPLA.potential\ N\ dθ\ dx\ U)\ ((N:ℝ≥0)*L)}
\leanline{8}{(kipldClockFriction\ N\ γ)\ (IPLA.potential\_gradient\_lipschitz\ hdiff\ hLip)}
\leanline{8}{(T/Real.sqrt\ N)\ (div\_nonneg\ hT\ (Real.sqrt\_nonneg\ \_))}
\end{leancode}
\leanlabel{How the Lean proof follows the argument}
\texttt{kinetic\_\allowbreak{}integral\_\allowbreak{}clock} transforms the full integral equation, including both drift coordinates. \texttt{clock\allowbreak{}Brownian\_\allowbreak{}is\allowbreak{}Pre\allowbreak{}Brownian} checks the transformed Brownian law. \texttt{kipld\allowbreak{}Brownian\allowbreak{}Position\allowbreak{}Law\_\allowbreak{}eq\_\allowbreak{}canonical} identifies the resulting position output measure with the unit-temperature kinetic output. This measure identity is the point at which the shared kinetic sampling theorem becomes applicable.
\end{leanbox}

\section{\texorpdfstring{Curvature and smoothness of the scaled particle potential}{Curvature and smoothness of the scaled particle potential}}


The gradients of the averaged potential are


\[\nabla_\theta\bar U_N=\frac1N\sum_i\nabla_\theta U(\theta,\sqrt Nz_i),
\qquad\nabla_{z_i}\bar U_N=\frac1{\sqrt N}\nabla_xU(\theta,\sqrt Nz_i).\]


Apply the lifted strong-convexity theorem from the earlier particle geometry report (see \cite[Section~3.1.3]{LangevinLean2026IPLA}). Equivalently, apply joint strong convexity to each pair
\((\theta,\sqrt Nz_i),(\theta',\sqrt Nz'_i)\), divide by \(N\), and sum.
The repeated parameter term and the scaled latent terms combine exactly:


\[\langle Z-Z',\nabla\bar U_N(Z)-\nabla\bar U_N(Z')\rangle
\ge\mu\left(\|\theta-\theta'\|^2+\sum_i\|z_i-z'_i\|^2\right).\]


For smoothness, Jensen bounds the squared norm of the averaged parameter
gradient by the average of squared norms. Adding the latent-gradient squares
then permits one joint \(L\)-Lipschitz estimate per particle, giving
\(\|\nabla\bar U_N(Z)-\nabla\bar U_N(Z')\|\le L\|Z-Z'\|\). This is the earlier lifted-gradient Lipschitz theorem divided by \(N\) (see \cite[Section~3.1.4]{LangevinLean2026IPLA}).




Thus the standard potential \(f_N=N\bar U_N\) has constants
\(m_N=N\mu\), \(M_N=NL\). Their ratio is \(L/\mu\), independent of
\(N\). These identities determine the dependence of the KIPLMC1 error bound on particle count.



The lifted sum potential has curvature and gradient Lipschitz constants multiplied by \(N\). Dividing the potential by \(N\) cancels both factors in physical time. As a result the synchronous rate depends on the original pair constants, while total coordinate dimension appears only through thermal and numerical moments.

\begin{leanbox}{Lean: \texttt{gradient\_\allowbreak{}averagePotential}, \texttt{averagePotential\_\allowbreak{}strongConvex}, \texttt{force\_\allowbreak{}lipschitz}}
\leansource{\texttt{gradient\_\allowbreak{}averagePotential} \quad KIPLMC/\allowbreak{}Geometry.lean, lines 30–39}
\begin{leancode}
\leanline{0}{theorem\ gradient\_averagePotential\ \{N\ dθ\ dx\ :\ ℕ\}\ \{U\ :\ IPLA.Joint\ dθ\ dx\ →\ ℝ\}}
\leanline{4}{(hdiff\ :\ Differentiable\ ℝ\ U)\ (z\ :\ State\ N\ dθ\ dx)\ :}
\leanline{4}{gradient\ (@averagePotential\ N\ dθ\ dx\ U)\ z\ =\ force\ U\ z}
\end{leancode}
\leansource{\texttt{averagePotential\_\allowbreak{}strongConvex} \quad KIPLMC/\allowbreak{}Geometry.lean, lines 48–58}
\begin{leancode}
\leanline{0}{theorem\ averagePotential\_strongConvex\ \{N\ dθ\ dx\ :\ ℕ\}\ (hN\ :\ 0\ <\ N)}
\leanline{4}{\{U\ :\ IPLA.Joint\ dθ\ dx\ →\ ℝ\}\ \{m\ :\ ℝ\}}
\leanline{4}{(hU\ :\ StrongConvexOn\ Set.univ\ m\ U)\ :}
\leanline{4}{StrongConvexOn\ Set.univ\ m\ (@averagePotential\ N\ dθ\ dx\ U)}
\end{leancode}
\leansource{\texttt{force\_\allowbreak{}lipschitz} \quad KIPLMC/\allowbreak{}Geometry.lean, lines 82–98}
\begin{leancode}
\leanline{0}{theorem\ force\_lipschitz\ \{N\ dθ\ dx\ :\ ℕ\}\ (hN\ :\ 0\ <\ N)}
\leanline{4}{\{U\ :\ IPLA.Joint\ dθ\ dx\ →\ ℝ\}\ \{L\ :\ ℝ≥0\}}
\leanline{4}{(hdiff\ :\ Differentiable\ ℝ\ U)\ (hLip\ :\ LipschitzWith\ L\ (gradient\ U))\ :}
\leanline{4}{LipschitzWith\ L\ (@force\ N\ dθ\ dx\ U)}
\end{leancode}
\leanlabel{How the Lean proof follows the argument}
\texttt{gradient\_\allowbreak{}average\allowbreak{}Potential} identifies the normalized gradient with the physical particle force. \texttt{average\allowbreak{}Potential\_\allowbreak{}strong\allowbreak{}Convex} and \texttt{force\_\allowbreak{}lipschitz} verify the curvature and Lipschitz inputs for that potential. Each theorem retains the same joint-potential assumptions, so no particle-dependent regularity constant is introduced when the kinetic theory is applied.
\end{leanbox}

\section{\texorpdfstring{The normalized powered-likelihood parameter marginal}{The normalized powered-likelihood parameter marginal}}


The rescaled kinetic Gibbs density is


\[\widetilde\Pi(dz,dv)\propto
\exp\left[-\sum_iU(\theta,\sqrt Nz_i)-\frac12\|v\|^2\right]dz\,dv.\]


Integrating the Gaussian velocity gives a constant. For each latent block,
substitute \(x_i=\sqrt Nz_i\). Its Jacobian is independent of \(\theta\),
so it is absorbed in normalization. The earlier exact parameter-marginal theorem identifies the normalized pushforward after the same coordinate change (see \cite[Section~3.9.9]{LangevinLean2026IPLA}). The remaining product separates:


\[\pi_\Theta(\theta)\propto
\prod_{i=1}^N\int e^{-U(\theta,x_i)}dx_i
=p_\theta(y)^N=e^{-N\kappa(\theta)},\quad
\kappa=-\log p_\theta(y).\]


At stationarity, more particles therefore sharpen the likelihood density
rather than changing its maximizer. The maximizer is unchanged by taking a positive integer power.
It relies on the parameter noise being present at the prescribed scale; a
noise-free parameter gradient flow has a different invariant structure.




The density computation identifies the candidate marginal. Normalization and finite second moments follow from the earlier lifted-target construction (see \cite[Section~3.5.1]{LangevinLean2026IPLA}, \cite[Section~3.5.1]{LangevinLean2026IPLA}). Invariance is a separate dynamical assertion, supplied by the kinetic phase-stationarity theorem after the clock change (see \cite[Section~3.58.3]{LangevinLean2026Kinetic}).



The latent variables become independent under the conditional Gibbs density for a fixed parameter. Fubini multiplies their identical integrals into \(\ell(\theta)^N\). Dividing by the finite positive full partition function gives the powered marginal, so the concentration argument applies to a normalized probability measure.

\begin{leanbox}{Lean: \texttt{gibbsPhaseTarget\_\allowbreak{}probability}, \texttt{gibbsPhaseTarget\_\allowbreak{}parameter}, \texttt{phaseFlow\_\allowbreak{}gibbs\_\allowbreak{}stationary}}
\leansource{\texttt{gibbsPhaseTarget\_\allowbreak{}probability} \quad KIPLMC/\allowbreak{}Target.lean, lines 35–41}
\begin{leancode}
\leanline{0}{theorem\ gibbsPhaseTarget\_probability\ \{N\ dθ\ dx\ :\ ℕ\}\ (hN\ :\ 0\ <\ N)}
\leanline{4}{(U\ :\ IPLA.Joint\ dθ\ dx\ →\ ℝ)\ (m\ :\ ℝ)\ (hm\ :\ 0\ <\ m)}
\leanline{4}{(hdiff\ :\ Differentiable\ ℝ\ U)\ (hsc\ :\ StrongConvexOn\ univ\ m\ U)\ :}
\leanline{4}{IsProbabilityMeasure\ (@gibbsPhaseTarget\ N\ dθ\ dx\ U)}
\end{leancode}
\leansource{\texttt{gibbsPhaseTarget\_\allowbreak{}parameter} \quad KIPLMC/\allowbreak{}Target.lean, lines 54–62}
\begin{leancode}
\leanline{0}{theorem\ gibbsPhaseTarget\_parameter\ \{N\ dθ\ dx\ :\ ℕ\}\ (hN\ :\ 0\ <\ N)}
\leanline{4}{(U\ :\ IPLA.Joint\ dθ\ dx\ →\ ℝ)\ (m\ :\ ℝ)\ (hm\ :\ 0\ <\ m)}
\leanline{4}{(hdiff\ :\ Differentiable\ ℝ\ U)\ (hsc\ :\ StrongConvexOn\ univ\ m\ U)\ :}
\leanline{4}{(@gibbsPhaseTarget\ N\ dθ\ dx\ U).map\ (fun\ z\ ↦\ IPLA.parameterCLM\ z.2)\ =}
\leanline{6}{IPLA.marginalTarget\ N\ U}
\end{leancode}
\leansource{\texttt{phaseFlow\_\allowbreak{}gibbs\_\allowbreak{}stationary} \quad KineticLangevin/\allowbreak{}ContinuousGibbs.lean, lines 89–118}
\begin{leancode}
\leanline{0}{theorem\ phaseFlow\_gibbs\_stationary\ (f\ :\ H\ ι\ →\ ℝ)\ (m\ :\ ℝ)\ (hm\ :\ 0\ <\ m)\ (M\ γ\ :\ ℝ≥0)}
\leanline{4}{(hf\ :\ Differentiable\ ℝ\ f)\ (hsc\ :\ StrongConvexOn\ univ\ m\ f)\ (hLip\ :\ LipschitzWith\ M\ (gradient\ f))}
\leanline{4}{(T\ :\ ℝ)\ (hT\ :\ 0\ ≤\ T)\ :}
\leanline{4}{letI\ :=\ ULA.isProbabilityMeasure\_gibbsMeasure\ f\ m\ hm\ hf\ hsc}
\leanline{4}{HasLaw\ (lpPhaseCoe\ ((ULA.gibbsMeasure\ f).prod\ (probability\ ι\ (stdGaussian\ (H\ ι))))}
\leanline{6}{(ProductContinuous.flow\ (probability\ ι\ (stdGaussian\ (H\ ι)))\ (ULA.gibbsMeasure\ f)}
\leanline{8}{(ULA.memLp\_id\_gibbsMeasure\ f\ m\ hm\ hf\ hsc)\ (velocity\ ι\ (stdGaussian\ (H\ ι))\ IsGaussian.memLp\_two\_id)}
\leanline{8}{(input\ ι\ (stdGaussian\ (H\ ι))\ γ)\ (continuous\_input\ ι\ (stdGaussian\ (H\ ι))\ γ)}
\leanline{8}{(gradient\ f)\ M\ γ\ hLip\ T\ hT\ T))}
\leanline{6}{(kineticGibbsTarget\ f)\ ((ULA.gibbsMeasure\ f).prod\ (probability\ ι\ (stdGaussian\ (H\ ι))))}
\end{leancode}
\leanlabel{How the Lean proof follows the argument}
\texttt{gibbs\allowbreak{}Phase\allowbreak{}Target\_\allowbreak{}probability} verifies normalization and \texttt{gibbs\allowbreak{}Phase\allowbreak{}Target\_\allowbreak{}parameter} computes the parameter marginal by integrating out velocities and latent coordinates. \texttt{phase\allowbreak{}Flow\_\allowbreak{}gibbs\_\allowbreak{}stationary} supplies invariance of the unit-temperature kinetic flow. The clock and law identities from the preceding part transfer this invariant distribution to the physical particle coordinates.
\end{leanbox}

\section{\texorpdfstring{The sharp parameter second moment about the maximizer}{The sharp parameter second moment about the maximizer}}


Strong Prékopa–Leindler preserves the parameter curvature when the latent
variable is integrated out: \(\kappa\) is \(\mu\)-strongly convex (see \cite[Section~3.12.3]{LangevinLean2026IPLA}). Thus its minimizer \(\theta_*\) is unique (see \cite[Section~3.12.4]{LangevinLean2026IPLA}), and the potential of \(\pi_\Theta\) has curvature \(N\mu\).




The sharp moment theorem in the earlier proximal report compares Gibbs partition integrals under dilation toward the minimizer (see \cite[Section~3.5.2]{LangevinLean2026Proximal}). It requires continuity and strong convexity of the marginal potential, without a differentiability assumption on that potential. Applied to \(N\kappa\), and then to the exact powered particle marginal (see \cite[Section~3.6.1]{LangevinLean2026Proximal}), it gives


\[W_2(\pi_\Theta,\delta_{\theta_*})
=\bigl(\mathbb E_{\pi_\Theta}\|\theta-\theta_*\|^2\bigr)^{1/2}
\le\sqrt{\frac{d_\theta}{N\mu}}.\]


The coefficient is sharp: an isotropic Gaussian with covariance
\((N\mu)^{-1}I\) attains equality. This explains why a proof with an extra
\(\sqrt2\) coefficient, although valid, would not reproduce Proposition 2
exactly.



The radius is a root-mean-square statement around the maximizer, rather than a covariance statement around the marginal mean. Its explicit coefficient is important for the final triangle inequality. No equality of the mean and maximizer is required to add this equilibrium contribution to the numerical sampling error.

\begin{leanbox}{Lean: \texttt{existsUnique\_\allowbreak{}mmle}, \texttt{parameter\_\allowbreak{}secondMoment\_\allowbreak{}le}, \texttt{parameter\_\allowbreak{}concentration}}
\leansource{\texttt{existsUnique\_\allowbreak{}mmle} \quad KIPLMC/\allowbreak{}Target.lean, lines 65–69}
\begin{leancode}
\leanline{0}{theorem\ existsUnique\_mmle\ \{dθ\ dx\ :\ ℕ\}\ (U\ :\ IPLA.Joint\ dθ\ dx\ →\ ℝ)}
\leanline{4}{(m\ :\ ℝ)\ (hm\ :\ 0\ <\ m)\ (hdiff\ :\ Differentiable\ ℝ\ U)}
\leanline{4}{(hsc\ :\ StrongConvexOn\ univ\ m\ U)\ :}
\leanline{4}{∃!\ θ,\ IPLA.IsMarginalMaximizer\ U\ θ}
\end{leancode}
\leansource{\texttt{parameter\_\allowbreak{}secondMoment\_\allowbreak{}le} \quad KIPLMC/\allowbreak{}Target.lean, lines 73–79}
\begin{leancode}
\leanline{0}{theorem\ parameter\_secondMoment\_le\ \{N\ dθ\ dx\ :\ ℕ\}\ (hN\ :\ 0\ <\ N)}
\leanline{4}{(U\ :\ IPLA.Joint\ dθ\ dx\ →\ ℝ)\ (m\ :\ ℝ)\ (hm\ :\ 0\ <\ m)}
\leanline{4}{(hdiff\ :\ Differentiable\ ℝ\ U)\ (hsc\ :\ StrongConvexOn\ univ\ m\ U)}
\leanline{4}{(θStar\ :\ EuclideanSpace\ ℝ\ (Fin\ dθ))\ (hmax\ :\ IPLA.IsMarginalMaximizer\ U\ θStar)\ :}
\leanline{4}{(∫\ θ,\ ‖θ\ -\ θStar‖\ \textasciicircum{}\ 2\ ∂IPLA.marginalTarget\ N\ U)\ ≤}
\leanline{6}{(dθ\ :\ ℝ)\ /\ ((N\ :\ ℝ)\ *\ m)}
\end{leancode}
\leansource{\texttt{parameter\_\allowbreak{}concentration} \quad KIPLMC/\allowbreak{}Target.lean, lines 82–89}
\begin{leancode}
\leanline{0}{theorem\ parameter\_concentration\ \{N\ dθ\ dx\ :\ ℕ\}\ (hN\ :\ 0\ <\ N)}
\leanline{4}{(U\ :\ IPLA.Joint\ dθ\ dx\ →\ ℝ)\ (m\ :\ ℝ)\ (hm\ :\ 0\ <\ m)}
\leanline{4}{(hdiff\ :\ Differentiable\ ℝ\ U)\ (hsc\ :\ StrongConvexOn\ univ\ m\ U)}
\leanline{4}{(θStar\ :\ EuclideanSpace\ ℝ\ (Fin\ dθ))\ (hmax\ :\ IPLA.IsMarginalMaximizer\ U\ θStar)\ :}
\leanline{4}{IPLA.projectedRMS\ (@IPLA.gibbsTarget\ N\ dθ\ dx\ U)\ IPLA.parameterCLM\ θStar\ ≤}
\leanline{6}{ENNReal.ofReal\ (Real.sqrt\ ((dθ\ :\ ℝ)\ /\ ((N\ :\ ℝ)\ *\ m)))}
\end{leancode}
\leanlabel{How the Lean proof follows the argument}
\texttt{exists\allowbreak{}Unique\_\allowbreak{}mmle} identifies the center. \texttt{parameter\_\allowbreak{}second\allowbreak{}Moment\_\allowbreak{}le} gives the dimension-sharp second moment about that center under the powered marginal; \texttt{parameter\_\allowbreak{}concentration} takes its square root and expresses the result in the projected root-mean-square form used by the transport triangle inequality.
\end{leanbox}

\section{\texorpdfstring{Transport to equilibrium and concentration about the maximizer}{Transport to equilibrium and concentration about the maximizer}}


The only coupling of a probability law with a point mass fixes the second
coordinate, so Wasserstein distance to a point is exactly root mean square
error. The triangle inequality gives


\[\bigl(\mathbb E\|\theta_n-\theta_*\|^2\bigr)^{1/2}
\le W_2(\operatorname{Law}(\theta_n),\pi_\Theta)
+\sqrt{\frac{d_\theta}{N\mu}}.\]


The first term is a numerical sampling error; the second is the equilibrium
statistical radius. Increasing \(n\) reduces the transient sampling error.
Decreasing the step reduces discretization bias. Increasing \(N\) reduces
the concentration term. These operations affect different pieces of the bound.




The map from full rescaled position, or full phase space, to its parameter
block is 1-Lipschitz. The earlier particle-geometry result gives operator norm at most one for this projection (see \cite[Section~3.1.6]{LangevinLean2026IPLA}). Applying it to any coupling cannot
increase its quadratic transport cost. Therefore a bound on the ordinary
kinetic sampler immediately controls the parameter sampling term once the
algorithm-law correspondence has been proved.



The projected root-mean-square transport inequality gives the next step (see \cite[Section~3.48.9]{LangevinLean2026ULASC}): for any coupling of the numerical position and an equilibrium position, project both to parameters and use Minkowski’s inequality around the maximizer. The projected discrepancy is bounded by the full scaled-state discrepancy because the projection is nonexpansive. Infimizing the latter cost yields the sampling term plus the equilibrium parameter radius.

\begin{leanbox}{Lean: \texttt{kipld\_\allowbreak{}parameter\_\allowbreak{}RMS\_\allowbreak{}bound}, \texttt{kiplmc1\_\allowbreak{}parameter\_\allowbreak{}bound}, \texttt{kiplmc2\_\allowbreak{}parameter\_\allowbreak{}bound}}
\leansource{\texttt{kipld\_\allowbreak{}parameter\_\allowbreak{}RMS\_\allowbreak{}bound} \quad KIPLMC/\allowbreak{}Continuous.lean, lines 248–269}
\begin{leancode}
\leanline{0}{theorem\ kipld\_parameter\_RMS\_bound\ \{N\ dθ\ dx\ :\ ℕ\}\ (hN\ :\ 0<N)}
\leanline{4}{(U\ :\ IPLA.Joint\ dθ\ dx\ →\ ℝ)\ (m\ :\ ℝ)\ (L\ γ\ :\ ℝ≥0)}
\leanline{4}{(hdiff\ :\ Differentiable\ ℝ\ U)\ (hsc\ :\ StrongConvexOn\ univ\ m\ U)}
\leanline{4}{(hLip\ :\ LipschitzWith\ L\ (gradient\ U))\ (hmL\ :\ m≤(L:ℝ))\ (hm\ :\ 0<m)}
\leanline{4}{(hγ\ :\ 0<(γ:ℝ))\ (hfriction\ :\ m+(L:ℝ)≤(γ:ℝ)\textasciicircum{}2)}
\leanline{4}{(ν\ :\ Measure\ (State\ N\ dθ\ dx))\ [IsProbabilityMeasure\ ν]}
\leanline{4}{(hν\ :\ MemLp\ (id\ :\ State\ N\ dθ\ dx\ →\ State\ N\ dθ\ dx)\ 2\ ν)\ (T\ :\ ℝ)\ (hT\ :\ 0≤T)}
\leanline{4}{(θStar\ :\ EuclideanSpace\ ℝ\ (Fin\ dθ))\ (hmax\ :\ IPLA.IsMarginalMaximizer\ U\ θStar)\ :}
\leanline{4}{Real.sqrt\ (∫\ z,\ ‖IPLA.parameterCLM\ z-θStar‖\textasciicircum{}2}
\leanline{6}{∂kipldBrownianPositionLaw\ hN\ U\ L\ γ\ hdiff\ hLip\ ν\ hν\ T\ hT)\ ≤}
\leanline{6}{Real.sqrt\ 2*Real.exp\ ((-m/(γ:ℝ))*T)*ULA.wasserstein2Real\ ν\ (IPLA.gibbsTarget\ U)+}
\leanline{6}{Real.sqrt\ ((dθ:ℝ)/((N:ℝ)*m))}
\end{leancode}
\leansource{\texttt{kiplmc1\_\allowbreak{}parameter\_\allowbreak{}bound} \quad KIPLMC/\allowbreak{}CanonicalExponential.lean, lines 146–169}
\begin{leancode}
\leanline{0}{theorem\ kiplmc1\_parameter\_bound\ \{N\ dθ\ dx\ :\ ℕ\}\ (hN\ :\ 0<N)}
\leanline{4}{(U\ :\ IPLA.Joint\ dθ\ dx→ℝ)\ (m\ :\ ℝ)\ (L\ γ\ :\ ℝ≥0)}
\leanline{4}{(hdiff\ :\ Differentiable\ ℝ\ U)\ (hsc\ :\ StrongConvexOn\ univ\ m\ U)}
\leanline{4}{(hLip\ :\ LipschitzWith\ L\ (gradient\ U))\ (hmL\ :\ m≤(L:ℝ))\ (hm\ :\ 0<m)}
\leanline{4}{(hγ\ :\ 0<(γ:ℝ))\ (hfriction\ :\ m+(L:ℝ)≤(γ:ℝ)\textasciicircum{}2)}
\leanline{4}{(h\ :\ ℝ≥0)\ (hh\ :\ 0<(h:ℝ))\ (hstep\ :\ 4*(γ:ℝ)*(L:ℝ)*(h:ℝ)≤m)}
\leanline{4}{(ν\ :\ Measure\ (State\ N\ dθ\ dx))\ [IsProbabilityMeasure\ ν]\ (hν\ :\ MemLp\ (id\ :\ State\ N\ dθ\ dx→State\ N\ dθ\ dx)\ 2\ ν)}
\leanline{4}{(θStar\ :\ EuclideanSpace\ ℝ\ (Fin\ dθ))\ (hmax\ :\ IPLA.IsMarginalMaximizer\ U\ θStar)\ (n\ :\ ℕ)\ :}
\leanline{4}{Real.sqrt\ (∫\ z,\ ‖IPLA.parameterCLM\ z-θStar‖\textasciicircum{}2}
\leanline{6}{∂(kiplmc1GaussianLaws\ U\ γ\ h\ ((velocityTarget\ N\ dθ\ dx).prod\ ν)\ n).map\ Prod.snd)\ ≤}
\leanline{6}{Real.sqrt\ 2*(1-3*m*(h:ℝ)/(4*(γ:ℝ)))\textasciicircum{}n*ULA.wasserstein2Real\ ν\ (IPLA.gibbsTarget\ U)+}
\leanline{6}{(L:ℝ)*(h:ℝ)*(Real.sqrt\ 2*Real.sqrt\ (((dθ:ℝ)+(N:ℝ)*(dx:ℝ))/(N:ℝ)))/m+}
\leanline{6}{Real.sqrt\ ((dθ:ℝ)/((N:ℝ)*m))}
\end{leancode}
\leansource{\texttt{kiplmc2\_\allowbreak{}parameter\_\allowbreak{}bound} \quad KIPLMC/\allowbreak{}OBABOBound.lean, lines 142–173}
\begin{leancode}
\leanline{0}{theorem\ kiplmc2\_parameter\_bound\ \{N\ dθ\ dx\ :\ ℕ\}\ (hN\ :\ 0<N)}
\leanline{4}{(U\ :\ IPLA.Joint\ dθ\ dx\ →\ ℝ)\ (m\ γ\ h\ :\ ℝ)\ (L\ :\ ℝ≥0)}
\leanline{4}{(hdiff\ :\ Differentiable\ ℝ\ U)\ (hsc\ :\ StrongConvexOn\ univ\ m\ U)}
\leanline{4}{(hLip\ :\ LipschitzWith\ L\ (gradient\ U))\ (hm\ :\ 0<m)\ (hmL\ :\ m≤(L\ :\ ℝ))}
\leanline{4}{(hγ\ :\ 0<γ)\ (hh\ :\ 0<h)\ (hfriction\ :\ 4*(L\ :\ ℝ)≤γ\textasciicircum{}2)\ (hstep\ :\ 33*γ\textasciicircum{}3*h≤m)}
\leanline{4}{(ν\ :\ Measure\ (State\ N\ dθ\ dx))\ [IsProbabilityMeasure\ ν]}
\leanline{4}{(hν\ :\ MemLp\ (id\ :\ State\ N\ dθ\ dx\ →\ State\ N\ dθ\ dx)\ 2\ ν)}
\leanline{4}{(θStar\ :\ EuclideanSpace\ ℝ\ (Fin\ dθ))\ (hmax\ :\ IPLA.IsMarginalMaximizer\ U\ θStar)\ (n\ :\ ℕ)\ :}
\leanline{4}{Real.sqrt\ (∫\ z,‖IPLA.parameterCLM\ z-θStar‖\textasciicircum{}2}
\leanline{6}{∂(kiplmc2Laws\ U\ γ\ h\ ((velocityTarget\ N\ dθ\ dx).prod\ ν)\ n).map\ Prod.snd)\ ≤}
\leanline{8}{Real.sqrt\ 2*(KineticLangevin.obaboContractionRate\ m\ γ\ h\ :\ ℝ)\textasciicircum{}n*}
\leanline{10}{ULA.wasserstein2Real\ ν\ (IPLA.gibbsTarget\ U)\ +}
\leanline{8}{Real.sqrt\ 2*((6*γ/m)*h*(Real.sqrt\ ((dθ\ :\ ℝ)+(N\ :\ ℝ)*(dx\ :\ ℝ))/Real.sqrt\ N)*}
\leanline{10}{(4*Real.sqrt\ (L\ :\ ℝ)+(L\ :\ ℝ)*h/2))\ +}
\leanline{8}{Real.sqrt\ ((dθ\ :\ ℝ)/((N\ :\ ℝ)*m))}
\end{leancode}
\leanlabel{How the Lean proof follows the argument}
\texttt{kipld\_\allowbreak{}parameter\_\allowbreak{}RMS\_\allowbreak{}bound}, \texttt{kiplmc1\_\allowbreak{}parameter\_\allowbreak{}bound} and \texttt{kiplmc2\_\allowbreak{}parameter\_\allowbreak{}bound} each apply the nonexpansive parameter projection and add \texttt{parameter\_\allowbreak{}concentration}. The coupling term is controlled in the complete scaled particle state before projection. Finiteness is proved before converting the extended transport and projected second moments to real square-root integrals.
\end{leanbox}

\section{\texorpdfstring{The KIPLMC1 exponential update and its change of coordinates}{The KIPLMC1 exponential update and its change of coordinates}}


Reuse \(\psi_0,\psi_1,\psi_2\) and the frozen-force exponential update from the earlier kinetic report (see \cite[Section~3.12.2]{LangevinLean2026Kinetic}). With fresh
correlated Gaussian pairs of covariance \(C\), KIPLMC1 updates


\[\theta^+=\theta+\psi_1(\eta)V^\theta-\psi_2(\eta)F^\theta
+\sqrt{2\gamma/N}\,\varepsilon^{0\prime},\]
\[V^{\theta+}=\psi_0(\eta)V^\theta-\psi_1(\eta)F^\theta
+\sqrt{2\gamma/N}\,\varepsilon^0.\]


Each latent block uses the same formulas with \(X^i,V^{x_i},F^{x_i}\)
and noise scale \(\sqrt{2\gamma}\). This is Algorithm 1.




Proposition B.1 proves the correspondence by changing variables in the
ordinary exponential step with \(\widetilde\gamma=\sqrt N\gamma\) and
\(\widetilde\eta=\eta/\sqrt N\). The coefficients obey


\[\widetilde\psi_0(\widetilde\eta)=\psi_0(\eta),\quad
\widetilde\psi_1(\widetilde\eta)=N^{-1/2}\psi_1(\eta),\quad
\widetilde\psi_2(\widetilde\eta)=N^{-1}\psi_2(\eta).\]


The velocity and position Gaussian integrals scale by \(N^{-1/4}\) and
\(N^{-3/4}\), respectively. Their covariance entries therefore scale by
\(N^{-1/2},N^{-1},N^{-3/2}\). Matching these laws, in addition to matching
the deterministic terms, permits reuse of the same kinetic convergence proof.



The parameter force contains a particle mean because the common parameter occurs in every potential summand. A latent force belongs to only one summand. Undoing its spatial scaling cancels the averaging factor, explaining why the physical latent update uses its individual gradient and ordinary thermal amplitude.

\begin{leanbox}{Lean: \texttt{kiplmc1Step}, \texttt{kiplmc1Step\_\allowbreak{}parameter\_\allowbreak{}position}, \texttt{kiplmc1Step\_\allowbreak{}particle\_\allowbreak{}velocity} and 1 more}
\leansource{\texttt{kiplmc1Step} \quad KIPLMC/\allowbreak{}Exponential.lean, lines 20–22}
\begin{leancode}
\leanline{0}{def\ kiplmc1Step\ \{N\ dθ\ dx\ :\ ℕ\}\ (U\ :\ IPLA.Joint\ dθ\ dx\ →\ ℝ)\ (γ\ h\ :\ ℝ)}
\leanline{4}{(z\ noise\ :\ State\ N\ dθ\ dx\ ×\ State\ N\ dθ\ dx)\ :=}
\leanline{2}{KineticLangevin.exponentialStep\ (force\ U)\ γ\ h\ z\ (kineticAmplitude\ N\ γ\ •\ noise)}
\end{leancode}
\leansource{\texttt{kiplmc1Step\_\allowbreak{}parameter\_\allowbreak{}position} \quad KIPLMC/\allowbreak{}Exponential.lean, lines 38–47}
\begin{leancode}
\leanline{0}{theorem\ kiplmc1Step\_parameter\_position\ \{N\ dθ\ dx\ :\ ℕ\}\ \{U\ :\ IPLA.Joint\ dθ\ dx\ →\ ℝ\}}
\leanline{4}{(hdiff\ :\ Differentiable\ ℝ\ U)\ (γ\ h\ :\ ℝ)}
\leanline{4}{(z\ noise\ :\ State\ N\ dθ\ dx\ ×\ State\ N\ dθ\ dx)\ (j\ :\ Fin\ dθ)\ :}
\leanline{4}{(kiplmc1Step\ U\ γ\ h\ z\ noise).2\ (Sum.inl\ j)\ =}
\leanline{6}{z.2\ (Sum.inl\ j)\ +\ KineticLangevin.psi\ γ\ 1\ h\ *\ z.1\ (Sum.inl\ j)\ -}
\leanline{6}{(KineticLangevin.psi\ γ\ 2\ h\ /\ N)\ *}
\leanline{8}{∑\ i\ :\ Fin\ N,\ gradient\ U\ (IPLA.joint\ i\ z.2)\ (Sum.inl\ j)\ +}
\leanline{6}{kineticAmplitude\ N\ γ\ *\ noise.2\ (Sum.inl\ j)}
\end{leancode}
\leansource{\texttt{kiplmc1Step\_\allowbreak{}particle\_\allowbreak{}velocity} \quad KIPLMC/\allowbreak{}Exponential.lean, lines 57–69}
\begin{leancode}
\leanline{0}{theorem\ kiplmc1Step\_particle\_velocity\ \{N\ dθ\ dx\ :\ ℕ\}\ (hN\ :\ 0\ <\ N)}
\leanline{4}{\{U\ :\ IPLA.Joint\ dθ\ dx\ →\ ℝ\}\ (hdiff\ :\ Differentiable\ ℝ\ U)\ (γ\ h\ :\ ℝ)}
\leanline{4}{(z\ noise\ :\ State\ N\ dθ\ dx\ ×\ State\ N\ dθ\ dx)\ (i\ :\ Fin\ N)\ (j\ :\ Fin\ dx)\ :}
\leanline{4}{Real.sqrt\ N\ *\ (kiplmc1Step\ U\ γ\ h\ z\ noise).1\ (Sum.inr\ (i,j))\ =}
\leanline{6}{KineticLangevin.psi\ γ\ 0\ h\ *\ (Real.sqrt\ N\ *\ z.1\ (Sum.inr\ (i,j)))\ -}
\leanline{6}{KineticLangevin.psi\ γ\ 1\ h\ *\ gradient\ U\ (IPLA.joint\ i\ z.2)\ (Sum.inr\ j)\ +}
\leanline{6}{Real.sqrt\ (2\ *\ γ)\ *\ noise.1\ (Sum.inr\ (i,j))}
\end{leancode}
\leansource{\texttt{kiplmc1GaussianPositionLaws\_\allowbreak{}eq\_\allowbreak{}canonical} \quad KIPLMC/\allowbreak{}CanonicalExponential.lean, lines 29–39}
\begin{leancode}
\leanline{0}{theorem\ kiplmc1GaussianPositionLaws\_eq\_canonical\ \{N\ dθ\ dx\ :\ ℕ\}\ (hN\ :\ 0<N)}
\leanline{4}{(U\ :\ IPLA.Joint\ dθ\ dx→ℝ)\ (L\ :\ ℝ≥0)}
\leanline{4}{(hdiff\ :\ Differentiable\ ℝ\ U)\ (hLip\ :\ LipschitzWith\ L\ (gradient\ U))}
\leanline{4}{(γ\ :\ ℝ)\ (h\ :\ ℝ≥0)\ (hγ\ :\ 0≤γ)}
\leanline{4}{(ν\ :\ Measure\ (State\ N\ dθ\ dx))\ [IsProbabilityMeasure\ ν]\ (n\ :\ ℕ)\ :}
\leanline{4}{(kiplmc1GaussianLaws\ U\ γ\ h\ ((velocityTarget\ N\ dθ\ dx).prod\ ν)\ n).map\ Prod.snd\ =}
\leanline{6}{GenericCanonicalCoupling.klmcLaw\ (@IPLA.potential\ N\ dθ\ dx\ U)\ (Real.sqrt\ N*γ)}
\leanline{8}{⟨(h:ℝ)/Real.sqrt\ N,div\_nonneg\ h.coe\_nonneg\ (Real.sqrt\_nonneg\ \_)⟩\ ν\ n}
\end{leancode}
\leanlabel{How the Lean proof follows the argument}
\texttt{kiplmc1\allowbreak{}Step} is the deterministic particle step. \texttt{kiplmc1\allowbreak{}Step\_\allowbreak{}parameter\_\allowbreak{}position} and \texttt{kiplmc1\allowbreak{}Step\_\allowbreak{}particle\_\allowbreak{}velocity} identify its coordinate blocks with the displayed physical equations. \texttt{kiplmc1\allowbreak{}Gaussian\allowbreak{}Position\allowbreak{}Laws\_\allowbreak{}eq\_\allowbreak{}canonical} transfers the complete Gaussian Markov law, including initialization and correlated innovations, to the canonical kinetic sampler at the changed clock.
\end{leanbox}

\section{\texorpdfstring{The physical-time KIPLMC1 error bound}{The physical-time KIPLMC1 error bound}}


Apply the constructed KLMC sampling theorem from the earlier kinetic report (see \cite[Section~3.57.19]{LangevinLean2026Kinetic}) with
\(m_N=N\mu\), \(M_N=NL\), \(h_N=\eta/\sqrt N\), and
\(\gamma_N=\sqrt N\gamma\). Then


\[\frac{m_Nh_N}{\gamma_N}=\frac{\mu\eta}{\gamma},\qquad
\frac{M_Nh_N}{m_N}\sqrt{2d_z}
=\sqrt2\frac L\mu\eta\sqrt{\frac{d_\theta}N+d_x}.\]


The linear friction and step restrictions reduce exactly to
\(\gamma\ge\sqrt{\mu+L}\) and \(\eta\le\mu/(4\gamma L)\).
The projected-RMS transport inequality (see \cite[Section~3.48.9]{LangevinLean2026ULASC}) and the sharp particle concentration radius (see \cite[Section~3.6.1]{LangevinLean2026Proximal}) then give


\[\mathcal E_n\le\sqrt2D_0
\left(1-\frac{3\mu\eta}{4\gamma}\right)^n
+\sqrt2\frac L\mu\eta\sqrt{\frac{d_\theta}N+d_x}
+\sqrt{\frac{d_\theta}{N\mu}}.\]


The initial position Wasserstein distance is a valid \(D_0\).
The paper instead upper-bounds it by the root mean square distance to an
independently sampled stationary position. Initialization must include mutually
independent Gaussian velocity blocks, independent of position, so that the
rescaled velocity is precisely the standard Gaussian product needed by the
shared theorem.




Choosing \(N\) of order \(d_\theta/\varepsilon^2\) and \(\eta\) of
order \(\varepsilon/\sqrt{d_x}\), subject to stability, bounds both persistent terms. The transient then needs order \(\sqrt{d_x}/\varepsilon\) steps,
up to logarithmic and curvature factors.



The total scaled dimension is \(D=d_\theta+Nd_x\), and thermal variance per coordinate is \(1/N\). Thus the stationary velocity radius is \(\sqrt{D/N}=\sqrt{d_\theta/N+d_x}\). This identity is the probabilistic meaning of the particle-factor cancellation in the first-order discretization term.

\begin{leanbox}{Lean: \texttt{kiplmc1\_\allowbreak{}parameter\_\allowbreak{}bound\_\allowbreak{}independent}, \texttt{kiplmc1\_\allowbreak{}accuracy}}
\leansource{\texttt{kiplmc1\_\allowbreak{}parameter\_\allowbreak{}bound\_\allowbreak{}independent} \quad KIPLMC/\allowbreak{}CanonicalExponential.lean, lines 208–232}
\begin{leancode}
\leanline{0}{theorem\ kiplmc1\_parameter\_bound\_independent\ \{N\ dθ\ dx\ :\ ℕ\}\ (hN\ :\ 0<N)}
\leanline{4}{(U\ :\ IPLA.Joint\ dθ\ dx→ℝ)\ (m\ :\ ℝ)\ (L\ γ\ :\ ℝ≥0)}
\leanline{4}{(hdiff\ :\ Differentiable\ ℝ\ U)\ (hsc\ :\ StrongConvexOn\ univ\ m\ U)}
\leanline{4}{(hLip\ :\ LipschitzWith\ L\ (gradient\ U))\ (hmL\ :\ m≤(L:ℝ))\ (hm\ :\ 0<m)}
\leanline{4}{(hγ\ :\ 0<(γ:ℝ))\ (hfriction\ :\ m+(L:ℝ)≤(γ:ℝ)\textasciicircum{}2)}
\leanline{4}{(h\ :\ ℝ≥0)\ (hh\ :\ 0<(h:ℝ))\ (hstep\ :\ 4*(γ:ℝ)*(L:ℝ)*(h:ℝ)≤m)}
\leanline{4}{(ν\ :\ Measure\ (State\ N\ dθ\ dx))\ [IsProbabilityMeasure\ ν]\ (hν\ :\ MemLp\ (id\ :\ State\ N\ dθ\ dx→State\ N\ dθ\ dx)\ 2\ ν)}
\leanline{4}{(θStar\ :\ EuclideanSpace\ ℝ\ (Fin\ dθ))\ (hmax\ :\ IPLA.IsMarginalMaximizer\ U\ θStar)\ (n\ :\ ℕ)\ :}
\leanline{4}{Real.sqrt\ (∫\ z,\ ‖IPLA.parameterCLM\ z-θStar‖\textasciicircum{}2}
\leanline{6}{∂(kiplmc1GaussianLaws\ U\ γ\ h\ ((velocityTarget\ N\ dθ\ dx).prod\ ν)\ n).map\ Prod.snd)\ ≤}
\leanline{6}{Real.sqrt\ 2*(1-3*m*(h:ℝ)/(4*(γ:ℝ)))\textasciicircum{}n*}
\leanline{8}{Real.sqrt\ (∫\ p\ :\ State\ N\ dθ\ dx×State\ N\ dθ\ dx,\ ‖p.1-p.2‖\textasciicircum{}2\ ∂ν.prod\ (IPLA.gibbsTarget\ U))+}
\leanline{6}{(L:ℝ)*(h:ℝ)*(Real.sqrt\ 2*Real.sqrt\ (((dθ:ℝ)+(N:ℝ)*(dx:ℝ))/(N:ℝ)))/m+}
\leanline{6}{Real.sqrt\ ((dθ:ℝ)/((N:ℝ)*m))}
\end{leancode}
\leansource{\texttt{kiplmc1\_\allowbreak{}accuracy} \quad KIPLMC/\allowbreak{}Accuracy.lean, lines 21–41}
\begin{leancode}
\leanline{0}{theorem\ kiplmc1\_accuracy\ \{N\ dθ\ dx\ :\ ℕ\}\ (hN\ :\ 0<N)}
\leanline{4}{(U\ :\ IPLA.Joint\ dθ\ dx→ℝ)\ (m\ :\ ℝ)\ (L\ γ\ :\ ℝ≥0)}
\leanline{4}{(hdiff\ :\ Differentiable\ ℝ\ U)\ (hsc\ :\ StrongConvexOn\ univ\ m\ U)}
\leanline{4}{(hLip\ :\ LipschitzWith\ L\ (gradient\ U))\ (hmL\ :\ m≤(L:ℝ))\ (hm\ :\ 0<m)}
\leanline{4}{(hγ\ :\ 0<(γ:ℝ))\ (hfriction\ :\ m+(L:ℝ)≤(γ:ℝ)\textasciicircum{}2)}
\leanline{4}{(h\ :\ ℝ≥0)\ (hh\ :\ 0<(h:ℝ))\ (hstep\ :\ 4*(γ:ℝ)*(L:ℝ)*(h:ℝ)≤m)}
\leanline{4}{(ν\ :\ Measure\ (State\ N\ dθ\ dx))\ [IsProbabilityMeasure\ ν]\ (hν\ :\ MemLp\ (id\ :\ State\ N\ dθ\ dx→State\ N\ dθ\ dx)\ 2\ ν)}
\leanline{4}{(θStar\ :\ EuclideanSpace\ ℝ\ (Fin\ dθ))\ (hmax\ :\ IPLA.IsMarginalMaximizer\ U\ θStar)}
\leanline{4}{(ε\ :\ ℝ)\ (hε\ :\ 0<ε)}
\leanline{4}{(hbias\ :\ (L:ℝ)*(h:ℝ)*(Real.sqrt\ 2*Real.sqrt\ (((dθ:ℝ)+(N:ℝ)*(dx:ℝ))/(N:ℝ)))/m≤ε/3)}
\leanline{4}{(hstat\ :\ Real.sqrt\ ((dθ:ℝ)/((N:ℝ)*m))≤ε/3)\ (n\ :\ ℕ)}
\leanline{4}{(hn\ :\ ULA.wasserstein2Real\ ν\ (IPLA.gibbsTarget\ U)=0\ ∨}
\leanline{6}{KineticLangevin.firstOrderIterationCount\ (3*m*(h:ℝ)/(4*(γ:ℝ)))}
\leanline{8}{(ULA.wasserstein2Real\ ν\ (IPLA.gibbsTarget\ U))\ ε≤n)\ :}
\leanline{4}{Real.sqrt\ (∫\ z,\ ‖IPLA.parameterCLM\ z-θStar‖\textasciicircum{}2}
\leanline{6}{∂(kiplmc1GaussianLaws\ U\ γ\ h\ ((velocityTarget\ N\ dθ\ dx).prod\ ν)\ n).map\ Prod.snd)≤ε}
\end{leancode}
\leanlabel{How the Lean proof follows the argument}
\texttt{kiplmc1\_\allowbreak{}parameter\_\allowbreak{}bound\_\allowbreak{}independent} uses the canonical first-order bound, the exact scaled rate and bias identities, parameter projection, and equilibrium concentration. The independent initial-position coupling supplies the displayed initial constant. \texttt{kiplmc1\_\allowbreak{}accuracy} then bounds the transient, discretization and concentration contributions separately and adds the three budgets.
\end{leanbox}

\section{\texorpdfstring{The five OBABO substeps}{The five OBABO substeps}}


KIPLMC2 uses the symmetric five-subflow composition established in the earlier kinetic report (see \cite[Section~3.20.1]{LangevinLean2026Kinetic}). Write
\(O\) for friction plus noise, \(B\) for the force kick, and \(A\) for free
transport. The composition is


\[O(\eta/2)\;B(\eta/2)\;A(\eta)\;B(\eta/2)\;O(\eta/2).\]


Let \(\delta=e^{-\gamma\eta/2}\),
\(s=\sqrt{1-\delta^2}\). For a standard velocity block, an \(O\) half-step
is \(v\mapsto\delta v+s\xi\); for the parameter block it is
\(V^\theta\mapsto\delta V^\theta+s\xi/\sqrt N\).
A \(B\) half-step subtracts \(\eta F/2\), and the full \(A\) step sends
\(x\mapsto x+\eta v\).




The second \(B\) kick must use the force at the newly computed position.
The last \(O\) step uses an independent Gaussian vector. Symmetry gives the
method its second-order splitting structure, while only first derivatives of
\(U\) are evaluated. This is an entirely different construction from the
Hessian-corrected KLMC2 in the first paper.



Each thermostat half step separately preserves the appropriate Gaussian velocity law, each force kick is a velocity translation, and each drift is a position translation. Their order determines which refreshed velocity moves position and where the second nonlinear force is evaluated.

\begin{leanbox}{Lean: \texttt{obaboStep\_\allowbreak{}expanded}, \texttt{kiplmc2Step}}
\leansource{\texttt{obaboStep\_\allowbreak{}expanded} \quad KineticLangevin/\allowbreak{}OBABO.lean, lines 43–48}
\begin{leancode}
\leanline{0}{theorem\ obaboStep\_expanded\ (force\ :\ E\ →\ E)\ (h\ δ\ σ\ :\ ℝ)\ (z\ :\ E\ ×\ E)}
\leanline{4}{(ξ\ ξ'\ :\ E)\ :}
\leanline{4}{obaboStep\ force\ h\ δ\ σ\ z\ ξ\ ξ'\ =}
\leanline{6}{let\ vhalf\ :=\ δ\ •\ z.1\ +\ σ\ •\ ξ\ -\ (h\ /\ 2)\ •\ force\ z.2}
\leanline{6}{let\ xnext\ :=\ z.2\ +\ h\ •\ vhalf}
\leanline{6}{(δ\ •\ (vhalf\ -\ (h\ /\ 2)\ •\ force\ xnext)\ +\ σ\ •\ ξ',\ xnext)}
\end{leancode}
\leansource{\texttt{kiplmc2Step} \quad KIPLMC/\allowbreak{}Splitting.lean, lines 20–23}
\begin{leancode}
\leanline{0}{def\ kiplmc2Step\ \{N\ dθ\ dx\ :\ ℕ\}\ (U\ :\ IPLA.Joint\ dθ\ dx\ →\ ℝ)\ (γ\ h\ :\ ℝ)}
\leanline{4}{(z\ :\ State\ N\ dθ\ dx\ ×\ State\ N\ dθ\ dx)\ (ξ\ ξ'\ :\ State\ N\ dθ\ dx)\ :=}
\leanline{2}{KineticLangevin.obaboStep\ (force\ U)\ h\ (KineticLangevin.obaboDamping\ γ\ h)}
\leanline{4}{(thermalAmplitude\ N\ γ\ h)\ z\ ξ\ ξ'}
\end{leancode}
\leanlabel{How the Lean proof follows the argument}
\texttt{obabo\allowbreak{}Step\_\allowbreak{}expanded} expands the thermostat, half-kick, drift, half-kick and thermostat composition. \texttt{kiplmc2\allowbreak{}Step} uses that same map with the particle force and thermal amplitude. Consequently the second force is evaluated at the updated position and the two Gaussian thermostat arguments stay distinct.
\end{leanbox}

\section{\texorpdfstring{The expanded OBABO update}{The expanded OBABO update}}


For one standard block with force \(g\), combine the first \(O\) and
\(B\) operations:


\[v_{1/2}=\delta v+s\xi-\frac\eta2g(x),
\qquad x^+=x+\eta v_{1/2}.\]


Apply the final force kick and OU refresh:


\[v^+=\delta\left(v_{1/2}-\frac\eta2g(x^+)\right)+s\xi'
=\delta^2v-\frac{\eta\delta}2\bigl(g(x)+g(x^+)\bigr)
+s(\delta\xi+\xi').\]


In the parameter block replace \(g\) by the average \(F^\theta\), and
both Gaussian amplitudes by \(s/\sqrt N\). The latent blocks use individual
forces and amplitude \(s\). These identities reproduce Algorithm 2 and
Proposition B.2.




The resulting law is determined by two independent standard Gaussian
vectors, independent of the current state. Proving the displayed algebra for
arbitrary supplied noise vectors verifies the deterministic update;
proving the stochastic theorem additionally constructs this Gaussian kernel,
its iterates and the required moment estimates.



The final combined noise is \(s(\delta\xi+\xi^{\prime})\), but position depends only on the first refresh through the half velocity. Therefore the full phase output has a nontrivial shared-noise structure even though the two input thermostat vectors are independent. The expanded formula retains this structure exactly.

\begin{leanbox}{Lean: \texttt{halfVelocity\_\allowbreak{}parameter}, \texttt{halfVelocity\_\allowbreak{}particle}, \texttt{kiplmc2Step\_\allowbreak{}velocity} and 1 more}
\leansource{\texttt{halfVelocity\_\allowbreak{}parameter} \quad KIPLMC/\allowbreak{}Splitting.lean, lines 44–53}
\begin{leancode}
\leanline{0}{theorem\ halfVelocity\_parameter\ \{N\ dθ\ dx\ :\ ℕ\}\ \{U\ :\ IPLA.Joint\ dθ\ dx\ →\ ℝ\}}
\leanline{4}{(hdiff\ :\ Differentiable\ ℝ\ U)\ (γ\ h\ :\ ℝ)}
\leanline{4}{(z\ :\ State\ N\ dθ\ dx\ ×\ State\ N\ dθ\ dx)\ (ξ\ :\ State\ N\ dθ\ dx)\ (j\ :\ Fin\ dθ)\ :}
\leanline{4}{halfVelocity\ U\ γ\ h\ z\ ξ\ (Sum.inl\ j)\ =}
\leanline{6}{KineticLangevin.obaboDamping\ γ\ h\ *\ z.1\ (Sum.inl\ j)\ +}
\leanline{6}{thermalAmplitude\ N\ γ\ h\ *\ ξ\ (Sum.inl\ j)\ -}
\leanline{6}{h\ /\ (2\ *\ N)\ *\ ∑\ i\ :\ Fin\ N,\ gradient\ U\ (IPLA.joint\ i\ z.2)\ (Sum.inl\ j)}
\end{leancode}
\leansource{\texttt{halfVelocity\_\allowbreak{}particle} \quad KIPLMC/\allowbreak{}Splitting.lean, lines 64–76}
\begin{leancode}
\leanline{0}{theorem\ halfVelocity\_particle\ \{N\ dθ\ dx\ :\ ℕ\}\ (hN\ :\ 0\ <\ N)}
\leanline{4}{\{U\ :\ IPLA.Joint\ dθ\ dx\ →\ ℝ\}\ (hdiff\ :\ Differentiable\ ℝ\ U)\ (γ\ h\ :\ ℝ)}
\leanline{4}{(z\ :\ State\ N\ dθ\ dx\ ×\ State\ N\ dθ\ dx)\ (ξ\ :\ State\ N\ dθ\ dx)}
\leanline{4}{(i\ :\ Fin\ N)\ (j\ :\ Fin\ dx)\ :}
\leanline{4}{Real.sqrt\ N\ *\ halfVelocity\ U\ γ\ h\ z\ ξ\ (Sum.inr\ (i,j))\ =}
\leanline{6}{KineticLangevin.obaboDamping\ γ\ h\ *\ (Real.sqrt\ N\ *\ z.1\ (Sum.inr\ (i,j)))\ +}
\leanline{6}{Real.sqrt\ (KineticLangevin.obaboVariance\ γ\ h)\ *\ ξ\ (Sum.inr\ (i,j))\ -}
\leanline{6}{(h\ /\ 2)\ *\ gradient\ U\ (IPLA.joint\ i\ z.2)\ (Sum.inr\ j)}
\end{leancode}
\leansource{\texttt{kiplmc2Step\_\allowbreak{}velocity} \quad KIPLMC/\allowbreak{}Splitting.lean, lines 37–41}
\begin{leancode}
\leanline{0}{theorem\ kiplmc2Step\_velocity\ \{N\ dθ\ dx\ :\ ℕ\}\ (U\ :\ IPLA.Joint\ dθ\ dx\ →\ ℝ)\ (γ\ h\ :\ ℝ)}
\leanline{4}{(z\ :\ State\ N\ dθ\ dx\ ×\ State\ N\ dθ\ dx)\ (ξ\ ξ'\ :\ State\ N\ dθ\ dx)\ :}
\leanline{4}{(kiplmc2Step\ U\ γ\ h\ z\ ξ\ ξ').1\ =\ KineticLangevin.obaboDamping\ γ\ h\ •}
\leanline{6}{(halfVelocity\ U\ γ\ h\ z\ ξ\ -\ (h\ /\ 2)\ •}
\leanline{8}{force\ U\ (z.2\ +\ h\ •\ halfVelocity\ U\ γ\ h\ z\ ξ))\ +\ thermalAmplitude\ N\ γ\ h\ •\ ξ'}
\end{leancode}
\leansource{\texttt{kiplmc2Laws\_\allowbreak{}probability} \quad KIPLMC/\allowbreak{}Laws.lean, lines 82–87}
\begin{leancode}
\leanline{0}{theorem\ kiplmc2Laws\_probability\ \{N\ dθ\ dx\ :\ ℕ\}\ \{U\ :\ IPLA.Joint\ dθ\ dx\ →\ ℝ\}}
\leanline{4}{\{L\ :\ ℝ≥0\}\ (hN\ :\ 0\ <\ N)\ (hdiff\ :\ Differentiable\ ℝ\ U)}
\leanline{4}{(hLip\ :\ LipschitzWith\ L\ (gradient\ U))\ (γ\ h\ :\ ℝ)}
\leanline{4}{(μ\ :\ Measure\ (State\ N\ dθ\ dx\ ×\ State\ N\ dθ\ dx))\ [IsProbabilityMeasure\ μ]\ (n\ :\ ℕ)\ :}
\leanline{4}{IsProbabilityMeasure\ (kiplmc2Laws\ U\ γ\ h\ μ\ n)}
\end{leancode}
\leanlabel{How the Lean proof follows the argument}
\texttt{half\allowbreak{}Velocity\_\allowbreak{}parameter} and \texttt{half\allowbreak{}Velocity\_\allowbreak{}particle} evaluate the scaled phase coordinates after the first thermostat and kick. \texttt{kiplmc2\allowbreak{}Step\_\allowbreak{}velocity} checks the second force and final thermostat in the expanded velocity formula. \texttt{kiplmc2\allowbreak{}Laws\_\allowbreak{}probability} verifies that iteration of this measurable map against the prescribed Gaussian source yields probability laws.
\end{leanbox}

\section{\texorpdfstring{Comparison through a discrete invariant measure}{Comparison through a discrete invariant measure}}


A finite-step OBABO kernel \(P_h\) generally has its own invariant law
\(\Pi_h\), different from the kinetic Gibbs law \(\Pi\). The cited
Monmarché analysis supplies phase-space contraction and invariant bias bounds
of the form


\[W_2(\rho P_h^n,\sigma P_h^n)\le A q^{n/2}W_2(\rho,\sigma),
\qquad b=W_2(\Pi_h,\Pi)\le B.\]


Apply the triangle inequality twice, using invariance of \(\Pi_h\):


\[W_2(\rho P_h^n,\Pi)
\le A q^{n/2}\bigl(W_2(\rho,\Pi)+B\bigr)+B.\]


This explains two different occurrences of the numerical bias: a persistent
\(B\) and a transient \(Aq^{n/2}B\). Dropping the latter would lose a term
present in the supplied paper's revised theorem.




These inequalities are substantial results about the OBABO kernel,
including existence of \(\Pi_h\). A theorem that takes them as hypotheses
checks the transport combination but does not prove the cited OBABO analysis.
The coverage audit distinguishes that conditional use from a completed
argument.



A second argument compares the chain with a transition preserving Gibbs and does not require a discrete invariant law. The earlier kinetic report proves that an exact \(\mathrm O\)–Hamiltonian–\(\mathrm O\) step preserves the continuous Gibbs target (see \cite[Section~3.120.3]{LangevinLean2026Kinetic}). Compare it with the OBABO update on the same initial state and the same Gaussian vectors. The Verlet local defect has \(L^2\) cost of order \(h^2\sqrt{d_z}\). The earlier kinetic perturbation theorem sums this cost directly over the chain (see \cite[Section~3.122.3]{LangevinLean2026Kinetic}). The resulting bound with the thermal dimension factor is stronger than the supplied coarse estimate, so it also implies every term retained in the printed theorem.

\begin{leanbox}{Lean: \texttt{exactObaboStep\_\allowbreak{}preserves\_\allowbreak{}target}, \texttt{obabo\_\allowbreak{}wasserstein\_\allowbreak{}bound}}
\leansource{\texttt{exactObaboStep\_\allowbreak{}preserves\_\allowbreak{}target} \quad KineticLangevin/\allowbreak{}ExactOBABO.lean, lines 54–85}
\begin{leancode}
\leanline{0}{theorem\ exactObaboStep\_preserves\_target}
\leanline{4}{(f\ :\ E\ →\ ℝ)\ (m\ :\ ℝ)\ (hm\ :\ 0<m)\ (M\ :\ ℝ≥0)}
\leanline{4}{(hf\ :\ Differentiable\ ℝ\ f)\ (hsc\ :\ StrongConvexOn\ univ\ m\ f)}
\leanline{4}{(hLip\ :\ LipschitzWith\ M\ (gradient\ f))\ (h\ :\ ℝ)\ (hh\ :\ 0≤h)}
\leanline{4}{(δ\ σ\ :\ ℝ)\ (hvariance\ :\ δ\textasciicircum{}2+σ\textasciicircum{}2=1)\ :}
\leanline{4}{((kineticGibbsTarget\ f).prod\ ((stdGaussian\ E).prod\ (stdGaussian\ E))).map}
\leanline{6}{(Function.uncurry\ (exactObaboStep\ f\ M\ hLip\ h\ hh\ δ\ σ))\ =\ kineticGibbsTarget\ f}
\end{leancode}
\leansource{\texttt{obabo\_\allowbreak{}wasserstein\_\allowbreak{}bound} \quad KineticLangevin/\allowbreak{}OBABOBias.lean, lines 134–184}
\begin{leancode}
\leanline{0}{theorem\ obabo\_wasserstein\_bound}
\leanline{4}{(f\ :\ EuclideanSpace\ ℝ\ ι\ →\ ℝ)\ (m\ γ\ h\ σ\ :\ ℝ)\ (M\ :\ ℝ≥0)}
\leanline{4}{(hf\ :\ Differentiable\ ℝ\ f)\ (hsc\ :\ StrongConvexOn\ univ\ m\ f)\ (hLip\ :\ LipschitzWith\ M\ (gradient\ f))}
\leanline{4}{(hm\ :\ 0\ <\ m)\ (hmM\ :\ m\ ≤\ (M\ :\ ℝ))\ (hγ\ :\ 0\ <\ γ)\ (hh\ :\ 0\ <\ h)}
\leanline{4}{(hfriction\ :\ 4\ *\ (M\ :\ ℝ)\ ≤\ γ\ \textasciicircum{}\ 2)\ (hstep\ :\ 33\ *\ γ\ \textasciicircum{}\ 3\ *\ h\ ≤\ m)}
\leanline{4}{(hvariance\ :\ obaboDamping\ γ\ h\ \textasciicircum{}\ 2\ +\ σ\ \textasciicircum{}\ 2\ =\ 1)}
\leanline{4}{(μ\ :\ Measure\ (EuclideanSpace\ ℝ\ ι\ ×\ EuclideanSpace\ ℝ\ ι))\ [IsProbabilityMeasure\ μ]}
\leanline{4}{(hfinite\ :\ ULA.wasserstein2\ (μ.map\ (weightedPhaseTransform\ (1\ /\ (M\ :\ ℝ))\ (1\ /\ γ)))}
\leanline{6}{((kineticGibbsTarget\ f).map\ (weightedPhaseTransform\ (1\ /\ (M\ :\ ℝ))\ (1\ /\ γ)))\ ≠\ ⊤)\ (n\ :\ ℕ)\ :}
\leanline{4}{ULA.wasserstein2Real}
\leanline{6}{((obaboLaws\ (gradient\ f)\ γ\ h\ σ}
\leanline{8}{((stdGaussian\ (EuclideanSpace\ ℝ\ ι)).prod\ (stdGaussian\ (EuclideanSpace\ ℝ\ ι)))\ μ\ n).map}
\leanline{10}{(weightedPhaseTransform\ (1\ /\ (M\ :\ ℝ))\ (1\ /\ γ)))}
\leanline{6}{((kineticGibbsTarget\ f).map\ (weightedPhaseTransform\ (1\ /\ (M\ :\ ℝ))\ (1\ /\ γ)))\ ≤}
\leanline{8}{(obaboContractionRate\ m\ γ\ h\ :\ ℝ)\ \textasciicircum{}\ n\ *}
\leanline{10}{ULA.wasserstein2Real\ (μ.map\ (weightedPhaseTransform\ (1\ /\ (M\ :\ ℝ))\ (1\ /\ γ)))}
\leanline{12}{((kineticGibbsTarget\ f).map\ (weightedPhaseTransform\ (1\ /\ (M\ :\ ℝ))\ (1\ /\ γ)))\ +}
\leanline{10}{(6\ *\ γ\ /\ m)\ *\ h\ *\ Real.sqrt\ (Fintype.card\ ι\ :\ ℝ)\ *\ (4\ *\ Real.sqrt\ (M\ :\ ℝ)\ +\ (M\ :\ ℝ)\ *\ h\ /\ 2)}
\end{leancode}
\leanlabel{How the Lean proof follows the argument}
\texttt{exact\allowbreak{}Obabo\allowbreak{}Step\_\allowbreak{}preserves\_\allowbreak{}target} proves invariance for the exact thermostat/\allowbreak{}Hamiltonian comparison step. \texttt{obabo\_\allowbreak{}wasserstein\_\allowbreak{}bound} applies the contractive random-map perturbation argument to the numerical map and that exact comparator. The common-noise local cost is accumulated along this comparison; no premise asserting a numerical invariant law is passed to the sampling theorem.
\end{leanbox}

\section{\texorpdfstring{Particle dependence of the OBABO coefficients}{Particle dependence of the OBABO coefficients}}


For comparison with the direct bound, the manuscript’s cited discrete-invariant-law analysis uses the following rescaled constants


\[A_N=\sqrt3\max\{\sqrt{NL},(NL)^{-1/2}\},\qquad
q=1-\frac{\mu\eta}{3\gamma}.\]


Under \(\gamma\ge2\sqrt L\) and
\(0<\eta\le\eta_0=\mu/(33\gamma^3)\), one has \(0<q<1\).
Write \(M_L=\max\{\sqrt L,L^{-1/2}\}\). Because \(N\ge1\),
\(A_N\le\sqrt{3N}M_L\). The detailed bias expression in the source
can be bounded by


\[B_N\le\sqrt3G_*\eta\sqrt{Nd_z},\]
\[G_*=\frac{6\gamma L}{\mu}M_L
\left(1+\frac{\eta_0L}{2\sqrt\mu}\right)
\left(1+e^{L\eta_0^2}
\left(\frac{\eta_0}{6}+\frac{\eta_0^2L}{24}\right)\right).\]


All factors increase with a nonnegative step, so evaluating at the fixed
upper limit \(\eta_0\) removes hidden step dependence. This makes \(G_*\)
a constant depending only on \(\mu,L,\gamma\).




The product \(A_NB_N\) is proportional to \(\eta N\sqrt{d_z}\), while
\(B_N\) itself is proportional to \(\eta\sqrt{Nd_z}\). The two terms have
different powers of \(N\), which is why the source tracks them separately
before introducing the implicit constant.



The clock changes both curvature constants and friction, while leaving the physical contraction ratio unchanged. Simplifying the bias only after substituting every factor avoids introducing spurious growing powers of \(N\). The direct position sampling theorem from the earlier kinetic report instead gives a bound with thermal dimension factor (see \cite[Section~3.123.4]{LangevinLean2026Kinetic}), which is stronger than the coarser printed bound.

\begin{leanbox}{Lean: \texttt{kiplmc2Laws\_\allowbreak{}unitTemperature}, \texttt{obaboBias\_\allowbreak{}unitTemperature}, \texttt{kiplmc2\_\allowbreak{}paper\_\allowbreak{}envelope}}
\leansource{\texttt{kiplmc2Laws\_\allowbreak{}unitTemperature} \quad KIPLMC/\allowbreak{}OBABORescaling.lean, lines 127–147}
\begin{leancode}
\leanline{0}{theorem\ kiplmc2Laws\_unitTemperature\ \{N\ dθ\ dx\ :\ ℕ\}\ (hN\ :\ 0<N)}
\leanline{4}{(U\ :\ IPLA.Joint\ dθ\ dx\ →\ ℝ)\ (L\ :\ ℝ≥0)}
\leanline{4}{(hdiff\ :\ Differentiable\ ℝ\ U)\ (hLip\ :\ LipschitzWith\ L\ (gradient\ U))\ (γ\ h\ :\ ℝ)}
\leanline{4}{(μ\ :\ Measure\ (State\ N\ dθ\ dx\ ×\ State\ N\ dθ\ dx))\ [IsProbabilityMeasure\ μ]\ (n\ :\ ℕ)\ :}
\leanline{4}{(kiplmc2Laws\ U\ γ\ h\ μ\ n).map\ (rescaleVelocity\ (Real.sqrt\ N))\ =}
\leanline{6}{KineticLangevin.obaboLaws\ (gradient\ (@IPLA.potential\ N\ dθ\ dx\ U))}
\leanline{8}{(Real.sqrt\ N*γ)\ (h/Real.sqrt\ N)}
\leanline{8}{(Real.sqrt\ (KineticLangevin.obaboVariance\ (Real.sqrt\ N*γ)\ (h/Real.sqrt\ N)))}
\leanline{8}{(obaboNoise\ N\ dθ\ dx)\ (μ.map\ (rescaleVelocity\ (Real.sqrt\ N)))\ n}
\end{leancode}
\leansource{\texttt{obaboBias\_\allowbreak{}unitTemperature} \quad KIPLMC/\allowbreak{}OBABOBound.lean, lines 69–74}
\begin{leancode}
\leanline{0}{theorem\ obaboBias\_unitTemperature\ (c\ m\ L\ γ\ h\ d\ :\ ℝ)\ (hc\ :\ 0<c)\ :}
\leanline{4}{(6*(c*γ)/(c\textasciicircum{}2*m))*(h/c)*Real.sqrt\ d*(4*Real.sqrt\ (c\textasciicircum{}2*L)+(c\textasciicircum{}2*L)*(h/c)/2)\ =}
\leanline{6}{(6*γ/m)*h*(Real.sqrt\ d/c)*(4*Real.sqrt\ L+L*h/2)}
\end{leancode}
\leansource{\texttt{kiplmc2\_\allowbreak{}paper\_\allowbreak{}envelope} \quad KIPLMC/\allowbreak{}OBABOPaperRate.lean, lines 52–108}
\begin{leancode}
\leanline{0}{theorem\ kiplmc2\_paper\_envelope\ (m\ L\ γ\ h\ N\ d\ dθ\ D\ error\ :\ ℝ)\ (n\ :\ ℕ)}
\leanline{4}{(hm\ :\ 0<m)\ (hL\ :\ 0≤L)\ (hγ\ :\ 0<γ)\ (hh\ :\ 0≤h)\ (hN\ :\ 1≤N)}
\leanline{4}{(hd\ :\ 0≤d)\ (hθ\ :\ 0≤dθ)\ (hD\ :\ 0≤D)\ (hstep\ :\ 33*γ\textasciicircum{}3*h≤m)}
\leanline{4}{(hbound\ :\ error≤Real.sqrt\ 2*(KineticLangevin.obaboContractionRate\ m\ γ\ h\ :\ ℝ)\textasciicircum{}n*D+}
\leanline{6}{Real.sqrt\ 2*((6*γ/m)*h*(Real.sqrt\ d/Real.sqrt\ N)*(4*Real.sqrt\ L+L*h/2))+}
\leanline{6}{Real.sqrt\ (dθ/(N*m)))\ :}
\leanline{4}{error≤kiplmc2PaperConstant\ m\ L\ γ*}
\leanline{6}{((Real.sqrt\ N*D+h*N*Real.sqrt\ d)*Real.exp\ (-m*h*(n\ :\ ℝ)/(6*γ))+}
\leanline{8}{h*Real.sqrt\ (N*d)+Real.sqrt\ (dθ/N))}
\end{leancode}
\leanlabel{How the Lean proof follows the argument}
\texttt{kiplmc2\allowbreak{}Laws\_\allowbreak{}unit\allowbreak{}Temperature} identifies the entire numerical law under velocity rescaling. \texttt{obabo\allowbreak{}Bias\_\allowbreak{}unit\allowbreak{}Temperature} cancels the powers of the clock factor in the bias. \texttt{kiplmc2\_\allowbreak{}paper\_\allowbreak{}envelope} bounds the remaining step-dependent coefficient uniformly on the admissible step range and enlarges the thermal dimension factor to the printed particle dependence.
\end{leanbox}

\section{\texorpdfstring{The KIPLMC2 bound with both transient terms}{The KIPLMC2 bound with both transient terms}}


Independent Gaussian initialization gives the same velocity marginal as
\(\widetilde\Pi\). Coupling these equal velocity blocks and using position
projection proves the exact identity


\[W_2(\rho_0,\widetilde\Pi)
=W_2(\nu,\widetilde\pi_Z)=D_0.\]


Insert the scaled \(A_N,B_N\) into the discrete equilibrium comparison,
project to \(\theta\), and add the concentration radius. With a constant
\(C\) depending only on \(\mu,L,\gamma\), the supplied manuscript obtains


\[\mathcal E_n\le C\left[
\bigl(\sqrt N D_0+\eta N\sqrt{d_\theta+Nd_x}\bigr)
e^{-\mu\eta n/(6\gamma)}
+\eta\sqrt{N(d_\theta+Nd_x)}
+\sqrt{d_\theta/N}\right].\]


The exponential follows from
\((1-\mu\eta/(3\gamma))^{n/2}\le e^{-\mu\eta n/(6\gamma)}\).
The paper gives an explicit common coefficient by taking the maximum of
\(\sqrt3M_L,\sqrt3G_*,3M_LG_*\) and absorbing the concentration factor
\(\mu^{-1/2}\).




This bound has worse particle dependence than the KIPLMC1 bound. The paper
regards that dependence as a limitation of the available estimate and compares
it with numerical behavior; the experiments do not themselves remove the
factor from the theorem.



The extra transient term is harmless for validity but matters when reproducing the displayed manuscript formula. It decays with the same rate as the initial transport discrepancy. The sharper invariant-comparator estimate from the earlier kinetic report bounds the chain directly (see \cite[Section~3.123.4]{LangevinLean2026Kinetic}) and implies this coarser statement by nonnegative termwise enlargements.

\begin{leanbox}{Lean: \texttt{kiplmc2\_\allowbreak{}parameter\_\allowbreak{}paper\_\allowbreak{}bound}, \texttt{kiplmc2\_\allowbreak{}accuracy}}
\leansource{\texttt{kiplmc2\_\allowbreak{}parameter\_\allowbreak{}paper\_\allowbreak{}bound} \quad KIPLMC/\allowbreak{}OBABOPaperRate.lean, lines 112–133}
\begin{leancode}
\leanline{0}{theorem\ kiplmc2\_parameter\_paper\_bound\ \{N\ dθ\ dx\ :\ ℕ\}\ (hN\ :\ 0<N)}
\leanline{4}{(U\ :\ IPLA.Joint\ dθ\ dx\ →\ ℝ)\ (m\ γ\ h\ :\ ℝ)\ (L\ :\ ℝ≥0)}
\leanline{4}{(hdiff\ :\ Differentiable\ ℝ\ U)\ (hsc\ :\ StrongConvexOn\ univ\ m\ U)}
\leanline{4}{(hLip\ :\ LipschitzWith\ L\ (gradient\ U))\ (hm\ :\ 0<m)\ (hmL\ :\ m≤(L\ :\ ℝ))}
\leanline{4}{(hγ\ :\ 0<γ)\ (hh\ :\ 0<h)\ (hfriction\ :\ 4*(L\ :\ ℝ)≤γ\textasciicircum{}2)\ (hstep\ :\ 33*γ\textasciicircum{}3*h≤m)}
\leanline{4}{(ν\ :\ Measure\ (State\ N\ dθ\ dx))\ [IsProbabilityMeasure\ ν]}
\leanline{4}{(hν\ :\ MemLp\ (id\ :\ State\ N\ dθ\ dx\ →\ State\ N\ dθ\ dx)\ 2\ ν)}
\leanline{4}{(θStar\ :\ EuclideanSpace\ ℝ\ (Fin\ dθ))\ (hmax\ :\ IPLA.IsMarginalMaximizer\ U\ θStar)\ (n\ :\ ℕ)\ :}
\leanline{4}{Real.sqrt\ (∫\ z,‖IPLA.parameterCLM\ z-θStar‖\textasciicircum{}2}
\leanline{6}{∂(kiplmc2Laws\ U\ γ\ h\ ((velocityTarget\ N\ dθ\ dx).prod\ ν)\ n).map\ Prod.snd)\ ≤}
\leanline{6}{kiplmc2PaperConstant\ m\ L\ γ*}
\leanline{8}{((Real.sqrt\ N*ULA.wasserstein2Real\ ν\ (IPLA.gibbsTarget\ U)+}
\leanline{10}{h*(N\ :\ ℝ)*Real.sqrt\ ((dθ\ :\ ℝ)+(N\ :\ ℝ)*(dx\ :\ ℝ)))*}
\leanline{12}{Real.exp\ (-m*h*(n\ :\ ℝ)/(6*γ))+}
\leanline{10}{h*Real.sqrt\ ((N\ :\ ℝ)*((dθ\ :\ ℝ)+(N\ :\ ℝ)*(dx\ :\ ℝ)))+}
\leanline{10}{Real.sqrt\ ((dθ\ :\ ℝ)/(N\ :\ ℝ)))}
\end{leancode}
\leansource{\texttt{kiplmc2\_\allowbreak{}accuracy} \quad KIPLMC/\allowbreak{}Accuracy.lean, lines 106–131}
\begin{leancode}
\leanline{0}{theorem\ kiplmc2\_accuracy\ \{N\ dθ\ dx\ :\ ℕ\}\ (hN\ :\ 0<N)}
\leanline{4}{(U\ :\ IPLA.Joint\ dθ\ dx→ℝ)\ (m\ γ\ h\ :\ ℝ)\ (L\ :\ ℝ≥0)}
\leanline{4}{(hdiff\ :\ Differentiable\ ℝ\ U)\ (hsc\ :\ StrongConvexOn\ univ\ m\ U)}
\leanline{4}{(hLip\ :\ LipschitzWith\ L\ (gradient\ U))\ (hm\ :\ 0<m)\ (hmL\ :\ m≤(L:ℝ))}
\leanline{4}{(hγ\ :\ 0<γ)\ (hh\ :\ 0<h)\ (hfriction\ :\ 4*(L:ℝ)≤γ\textasciicircum{}2)\ (hstep\ :\ 33*γ\textasciicircum{}3*h≤m)}
\leanline{4}{(ν\ :\ Measure\ (State\ N\ dθ\ dx))\ [IsProbabilityMeasure\ ν]}
\leanline{4}{(hν\ :\ MemLp\ (id\ :\ State\ N\ dθ\ dx→State\ N\ dθ\ dx)\ 2\ ν)}
\leanline{4}{(θStar\ :\ EuclideanSpace\ ℝ\ (Fin\ dθ))\ (hmax\ :\ IPLA.IsMarginalMaximizer\ U\ θStar)}
\leanline{4}{(ε\ :\ ℝ)\ (hε\ :\ 0<ε)}
\leanline{4}{(hbias\ :\ kiplmc2PaperConstant\ m\ L\ γ*h*Real.sqrt\ ((N:ℝ)*((dθ:ℝ)+(N:ℝ)*(dx:ℝ)))≤ε/3)}
\leanline{4}{(hstat\ :\ kiplmc2PaperConstant\ m\ L\ γ*Real.sqrt\ ((dθ:ℝ)/(N:ℝ))≤ε/3)\ (n\ :\ ℕ)}
\leanline{4}{(hn\ :\ let\ D\ :=\ kiplmc2PaperConstant\ m\ L\ γ*(Real.sqrt\ N*ULA.wasserstein2Real\ ν\ (IPLA.gibbsTarget\ U)+}
\leanline{10}{h*(N:ℝ)*Real.sqrt\ ((dθ:ℝ)+(N:ℝ)*(dx:ℝ)))}
\leanline{6}{D=0\ ∨\ KineticLangevin.firstOrderIterationCount\ (m*h/(6*γ))\ D\ ε≤n)\ :}
\leanline{4}{Real.sqrt\ (∫\ z,‖IPLA.parameterCLM\ z-θStar‖\textasciicircum{}2}
\leanline{6}{∂(kiplmc2Laws\ U\ γ\ h\ ((velocityTarget\ N\ dθ\ dx).prod\ ν)\ n).map\ Prod.snd)≤ε}
\end{leancode}
\leanlabel{How the Lean proof follows the argument}
\texttt{kiplmc2\_\allowbreak{}parameter\_\allowbreak{}paper\_\allowbreak{}bound} follows from the stronger direct position sampling estimate, the uniform envelope and the sharp equilibrium concentration bound. Both transient terms of the printed theorem are retained as nonnegative upper bounds. \texttt{kiplmc2\_\allowbreak{}accuracy} then applies the chosen positive step, particle-count lower bound and rounded iteration requirement to the corresponding separate budgets.
\end{leanbox}

\section{\texorpdfstring{Iteration counts and the scope of the experimental examples}{Iteration counts and the scope of the experimental examples}}


For KIPLMC2, choosing \(N\asymp d_\theta/\varepsilon^2\) gives


\[\sqrt{N(d_\theta+Nd_x)}
\asymp\frac{d_\theta}{\varepsilon}
\sqrt{1+\frac{d_x}{\varepsilon^2}}.\]


To make the persistent numerical error of order \(\varepsilon\), take
the minimum of the stability threshold and
\(\varepsilon/\sqrt{N(d_\theta+Nd_x)}\). In the small-error regime this
is a step of order \(\varepsilon^3/(d_\theta\sqrt{d_x})\). Both transient
terms can then be controlled after
\(\widetilde O(d_\theta\sqrt{d_x}/\varepsilon^3)\) steps, provided the
initial transport error grows at most polynomially in particle number and
dimension so that its logarithm remains in the suppressed factor.




The manuscript also studies logistic regression, Wisconsin cancer data, a
Bayesian neural network, a multi-well model and runtime comparisons. Its signed
weighted area between error curves is
\(\sum_{i=1}^M2i(c^{\mathrm{IPLA}}(i)-c^{\mathrm{KIPLMC2}}(i))/(M(M+1))\).
Positive values favor KIPLMC2 by this metric. Its predictive log score averages
the log probability assigned to each correct test label.




These are empirical comparisons with concrete data and numerical settings.
The strong-convexity convergence theorems, their asymptotic complexity
corollaries, the algorithm definitions and the experimental metrics are
different kinds of content. The logistic example is globally smooth and can be strictly convex under a full-span feature condition, while common-translation directions still prevent a positive uniform strong-convexity constant without additional regularization.



The statistical choice of particle count first fixes the dimension-dependent numerical coefficient. The discretization-error bound then fixes a step, and the contraction rate fixes the iteration count. Empirical advantages at particular model settings remain separate from these global sufficient conditions, especially when a model lacks uniform joint strong convexity.

\begin{leanbox}{Lean: \texttt{accuracyParticleCount\_\allowbreak{}spec}, \texttt{accuracyStep\_\allowbreak{}spec}, \texttt{complexity\_\allowbreak{}dimension\_\allowbreak{}identity} and 1 more}
\leansource{\texttt{accuracyParticleCount\_\allowbreak{}spec} \quad KIPLMC/\allowbreak{}Accuracy.lean, lines 63–79}
\begin{leancode}
\leanline{0}{theorem\ accuracyParticleCount\_spec\ (C\ d\ ε\ :\ ℝ)\ (hC\ :\ 0≤C)\ (hd\ :\ 0≤d)\ (hε\ :\ 0<ε)\ :}
\leanline{4}{0<accuracyParticleCount\ C\ d\ ε\ ∧}
\leanline{4}{C*Real.sqrt\ (d/(accuracyParticleCount\ C\ d\ ε:ℝ))≤ε/3}
\end{leancode}
\leansource{\texttt{accuracyStep\_\allowbreak{}spec} \quad KIPLMC/\allowbreak{}Accuracy.lean, lines 86–90}
\begin{leancode}
\leanline{0}{theorem\ accuracyStep\_spec\ (A\ B\ ε\ :\ ℝ)\ (hA\ :\ 0<A)\ (hB\ :\ 0<B)\ (hε\ :\ 0<ε)\ :}
\leanline{4}{0<accuracyStep\ A\ B\ ε\ ∧\ accuracyStep\ A\ B\ ε≤A\ ∧\ accuracyStep\ A\ B\ ε*B≤ε/3}
\end{leancode}
\leansource{\texttt{complexity\_\allowbreak{}dimension\_\allowbreak{}identity} \quad KIPLMC/\allowbreak{}Accuracy.lean, lines 93–103}
\begin{leancode}
\leanline{0}{theorem\ complexity\_dimension\_identity\ (dθ\ dx\ ε\ :\ ℝ)\ (hθ\ :\ 0≤dθ)\ (hx\ :\ 0≤dx)\ (hε\ :\ 0<ε)\ :}
\leanline{4}{Real.sqrt\ ((dθ/ε\textasciicircum{}2)*(dθ+(dθ/ε\textasciicircum{}2)*dx))=}
\leanline{6}{(dθ/ε)*Real.sqrt\ (1+dx/ε\textasciicircum{}2)}
\end{leancode}
\leansource{\texttt{kiplmc2\_\allowbreak{}accuracy} \quad KIPLMC/\allowbreak{}Accuracy.lean, lines 106–131}
\begin{leancode}
\leanline{0}{theorem\ kiplmc2\_accuracy\ \{N\ dθ\ dx\ :\ ℕ\}\ (hN\ :\ 0<N)}
\leanline{4}{(U\ :\ IPLA.Joint\ dθ\ dx→ℝ)\ (m\ γ\ h\ :\ ℝ)\ (L\ :\ ℝ≥0)}
\leanline{4}{(hdiff\ :\ Differentiable\ ℝ\ U)\ (hsc\ :\ StrongConvexOn\ univ\ m\ U)}
\leanline{4}{(hLip\ :\ LipschitzWith\ L\ (gradient\ U))\ (hm\ :\ 0<m)\ (hmL\ :\ m≤(L:ℝ))}
\leanline{4}{(hγ\ :\ 0<γ)\ (hh\ :\ 0<h)\ (hfriction\ :\ 4*(L:ℝ)≤γ\textasciicircum{}2)\ (hstep\ :\ 33*γ\textasciicircum{}3*h≤m)}
\leanline{4}{(ν\ :\ Measure\ (State\ N\ dθ\ dx))\ [IsProbabilityMeasure\ ν]}
\leanline{4}{(hν\ :\ MemLp\ (id\ :\ State\ N\ dθ\ dx→State\ N\ dθ\ dx)\ 2\ ν)}
\leanline{4}{(θStar\ :\ EuclideanSpace\ ℝ\ (Fin\ dθ))\ (hmax\ :\ IPLA.IsMarginalMaximizer\ U\ θStar)}
\leanline{4}{(ε\ :\ ℝ)\ (hε\ :\ 0<ε)}
\leanline{4}{(hbias\ :\ kiplmc2PaperConstant\ m\ L\ γ*h*Real.sqrt\ ((N:ℝ)*((dθ:ℝ)+(N:ℝ)*(dx:ℝ)))≤ε/3)}
\leanline{4}{(hstat\ :\ kiplmc2PaperConstant\ m\ L\ γ*Real.sqrt\ ((dθ:ℝ)/(N:ℝ))≤ε/3)\ (n\ :\ ℕ)}
\leanline{4}{(hn\ :\ let\ D\ :=\ kiplmc2PaperConstant\ m\ L\ γ*(Real.sqrt\ N*ULA.wasserstein2Real\ ν\ (IPLA.gibbsTarget\ U)+}
\leanline{10}{h*(N:ℝ)*Real.sqrt\ ((dθ:ℝ)+(N:ℝ)*(dx:ℝ)))}
\leanline{6}{D=0\ ∨\ KineticLangevin.firstOrderIterationCount\ (m*h/(6*γ))\ D\ ε≤n)\ :}
\leanline{4}{Real.sqrt\ (∫\ z,‖IPLA.parameterCLM\ z-θStar‖\textasciicircum{}2}
\leanline{6}{∂(kiplmc2Laws\ U\ γ\ h\ ((velocityTarget\ N\ dθ\ dx).prod\ ν)\ n).map\ Prod.snd)≤ε}
\end{leancode}
\leanlabel{How the Lean proof follows the argument}
\texttt{accuracy\allowbreak{}Particle\allowbreak{}Count\_\allowbreak{}spec} proves the statistical radius bound for the rounded natural particle count. \texttt{accuracy\allowbreak{}Step\_\allowbreak{}spec} verifies positivity and the step upper bounds. \texttt{complexity\_\allowbreak{}dimension\_\allowbreak{}identity} checks the dimension algebra in the resulting count, while \texttt{kiplmc2\_\allowbreak{}accuracy} supplies the finite error guarantee. Timing traces and empirical predictive performance are not conclusions of these declarations.
\end{leanbox}

\begin{thebibliography}{ULA-LSI}
\bibitem[ULA-SC]{LangevinLean2026ULASC}
LangevinLean Project. \emph{Unadjusted Langevin under Strong Convexity}. LL-01, 2026. Draft report, content revision d50c0f64815b. \url{https://odakyildiz.github.io/LangevinLean/docs/pdf/UnadjustedLangevin/FullProof.pdf}.
\bibitem[IPLA-SC]{LangevinLean2026IPLA}
LangevinLean Project. \emph{Interacting Particle Langevin under Strong Convexity}. LL-03, 2026. Draft report, content revision bfbb88fd1d79. \url{https://odakyildiz.github.io/LangevinLean/docs/pdf/IPLA/FullProof.pdf}.
\bibitem[Proximal]{LangevinLean2026Proximal}
LangevinLean Project. \emph{Proximal particle Langevin}. LL-05, 2026. Draft report, content revision f201d59ac57b. \url{https://odakyildiz.github.io/LangevinLean/docs/pdf/ProximalLangevin/FullProof.pdf}.
\bibitem[Kinetic]{LangevinLean2026Kinetic}
LangevinLean Project. \emph{Kinetic Langevin Monte Carlo}. LL-06, 2026. Draft report, content revision 74e68143120e. \url{https://odakyildiz.github.io/LangevinLean/docs/pdf/KineticLangevin/FullProof.pdf}.
\end{thebibliography}

\end{document}
